% ch05 第5章习题解答（21 题）
\chapter{相互作用费米子系统}

% ---------------------------------------------------------------
\section*{问题 5.1.1}
\begin{problem}
对式 (5.1.1) 的 X 射线模型（芯能级能量 $-E_C$，耦合 $\Gamma$），计算吸收率谱 $A_{abs}(\omega)$，即芯电子被光子激发到导带的过程的速率谱。\emph{（译注：原书题面排印为 $A_{obs}$，按上下文应为吸收率 $A_{abs}$。）}
\end{problem}
\solution

吸收过程是发射过程的逆过程：光子被吸收，芯电子（能量 $-E_C$，以费米能为零点）被激发到导带中能量为 $\xi_{\pmb{k}}$ 的未占据态。能量守恒要求
\begin{equation*}
\omega=(-\xi_{\pmb{k}})-(-E_C)\quad\Longrightarrow\quad \xi_{\pmb{k}}=\omega-E_C .
\end{equation*}
与发射率的推导完全平行（把隧穿算符换成 $I^{\dagger}=c^{\dagger}(0)C$，即从芯态取走一个电子、在导带放一个电子），由线性响应理论得
\begin{equation*}
A_{abs}(\omega)=-2\pi\Gamma^{2}\int\mathrm{d}\nu\,\big(A_{C-}^{0}(\nu)\,A_{c+}^{0}(\nu-\omega)-A_{C+}^{0}(\nu)\,A_{c-}^{0}(\nu-\omega)\big).
\end{equation*}
吸收前芯态被完全填满，而非平衡谱函数为
\begin{equation*}
A_{C-}^{0}(\omega)=\delta(\omega+E_C),\qquad A_{C+}^{0}(\omega)=0 ,
\end{equation*}
故第二项为零，第一项中 $\delta(\nu+E_C)$ 把 $\nu$ 固定在 $\nu=\omega-E_C$：
\begin{equation*}
A_{abs}(\omega)=-2\pi\Gamma^{2}A_{c+}^{0}(\omega-E_C)
=2\pi\Gamma^{2}\big(1-n_F(\omega-E_C)\big)\,N(E_F-E_C+\omega).
\end{equation*}
即
\begin{equation*}
\boxed{\;A_{abs}(\omega)=2\pi\Gamma^{2}\big(1-n_F(\omega-E_C)\big)\,N(E_F-E_C+\omega)\;}
\end{equation*}
零温下 $A_{abs}(\omega)=2\pi\Gamma^{2}N(E_F-E_C+\omega)\,\Theta(\omega-E_C)$：与发射率（在 $\omega=E_C$ 处向下的阶梯）恰好互为镜像——吸收只在阈值 $\omega=E_C$ 之上开始，且直接测量费米面以上的态密度。物理上，$\omega<E_C$ 时吸收一个光子不足以把芯电子送上费米面，故无吸收。

% ---------------------------------------------------------------
\section*{问题 5.1.2}
\begin{problem}
对流体动力学哈密顿量 (5.1.5)（$A_{\pmb{k}}=m/\mathcal{V}\rho_0\pmb{k}^2$，$B_{\pmb{k}}=1/\mathcal{V}\chi$），把 $\alpha_{\pmb{k}}$ 用 $\rho_{\pmb{k}}$ 与 $\dot{\rho}_{\pmb{k}}$ 表示，并证明 $\alpha_{\pmb{k}}$ 携带确定的动量。
\end{problem}
\solution

以 $\rho_{\pmb{k}}$ 为正则坐标，由式 (5.1.4) 得正则动量
\begin{equation*}
\pi_{\pmb{k}}=\frac{\partial L}{\partial\dot{\rho}_{\pmb{k}}}
=\frac{m}{\mathcal{V}\rho_0\pmb{k}^{2}}\,\dot{\rho}_{-\pmb{k}}=A_{\pmb{k}}\dot{\rho}_{-\pmb{k}} .
\end{equation*}
仿照式 (3.3.17) 引入
\begin{equation*}
\alpha_{\pmb{k}}=u_{\pmb{k}}\rho_{\pmb{k}}+\mathrm{i}v_{\pmb{k}}\pi_{-\pmb{k}},
\qquad
u_{\pmb{k}}=\tfrac{1}{\sqrt2}(A_{\pmb{k}}B_{\pmb{k}})^{1/4},
\quad
v_{\pmb{k}}=\tfrac{1}{\sqrt2}(A_{\pmb{k}}B_{\pmb{k}})^{-1/4},
\end{equation*}
并注意 $\Omega_{\pmb{k}}=\sqrt{B_{\pmb{k}}/A_{\pmb{k}}}$，$A_{\pmb{k}}v_{\pmb{k}}=u_{\pmb{k}}/\Omega_{\pmb{k}}$（直接验证：$\tfrac1{\sqrt2}A^{3/4}B^{-1/4}=\tfrac1{\sqrt2}(AB)^{1/4}\sqrt{A/B}$），于是
\begin{equation*}
\boxed{\;\alpha_{\pmb{k}}=u_{\pmb{k}}\Big(\rho_{\pmb{k}}+\mathrm{i}\frac{\dot{\rho}_{\pmb{k}}}{\Omega_{\pmb{k}}}\Big),
\qquad
u_{\pmb{k}}=\Big(\frac{m}{4\mathcal{V}^{2}\rho_0\chi\,\pmb{k}^{2}}\Big)^{1/4},
\quad
\Omega_{\pmb{k}}=\sqrt{\frac{\rho_0}{m\chi}}\,|\pmb{k}|=v|\pmb{k}|\;}
\end{equation*}
（这正是正文用来写 $\delta H=\mathcal{V}^{-1}\sum_{\pmb{k}}V_{-\pmb{k}}(\alpha_{\pmb{k}}+\alpha_{\pmb{k}}^{\dagger})/2u_{\pmb{k}}$ 的关系式 $\rho_{\pmb{k}}=(\alpha_{\pmb{k}}+\alpha_{\pmb{k}}^{\dagger})/2u_{\pmb{k}}$ 的逆。）

动量：总动量 $\pmb{P}$ 生成空间平移，$U(\pmb{a})\rho(\pmb{x})U^{\dagger}(\pmb{a})=\rho(\pmb{x}+\pmb{a})$，故对 $\rho_{\pmb{k}}=\int\mathrm{d}^dx\,\rho(\pmb{x})\mathrm{e}^{-\mathrm{i}\pmb{k}\cdot\pmb{x}}$ 有
\begin{equation*}
[\pmb{P},\rho_{\pmb{k}}]=\pmb{k}\,\rho_{\pmb{k}} .
\end{equation*}
$\dot{\rho}_{\pmb{k}}$ 是 $\rho_{\pmb{k}}$ 的时间导数，动量本征值不变：$[\pmb{P},\dot{\rho}_{\pmb{k}}]=\pmb{k}\,\dot{\rho}_{\pmb{k}}$。因此
\begin{equation*}
[\pmb{P},\alpha_{\pmb{k}}]=\pmb{k}\,\alpha_{\pmb{k}} ,
\end{equation*}
即 $\alpha_{\pmb{k}}$ 作用在动量本征态上使其动量增加 $\pmb{k}$：\textbf{$\alpha_{\pmb{k}}$ 携带确定的动量 $\pmb{k}$}（$\alpha_{\pmb{k}}^{\dagger}$ 携带 $-\pmb{k}$）。这正是准粒子（声子）动量量子化的来源。

% ---------------------------------------------------------------
\section*{问题 5.1.3}
\begin{problem}
证明一维流体动力学模型（固定粒子数）与（无自旋）费米子系统（固定费米子数）的低能激发谱完全相同：(a) 能量 $nk_0v_F$、动量 $nk_0$ 的态有 $p_n$ 个（分割数，$k_0=2\pi/L$）；(b) 动量 $-nk_0$ 者亦然；(c) 能量 $nk_0v_F$、动量 $(n-2m)k_0$ 的态有 $p_{n-m}p_m$ 个。
\end{problem}
\solution

\textbf{流体动力学一侧。}哈密顿量为 $H=\sum_{\pmb{k}\neq0}v|\pmb{k}|\,\alpha_{\pmb{k}}^{\dagger}\alpha_{\pmb{k}}$（$v=v_F$），$\pmb{k}$ 取 $k_0$ 的整数倍。“固定粒子数”意味着弃去 $\pmb{k}=0$ 模（它改变总密度）。设激发态在动量 $jk_0$（$j=\pm1,\pm2,\dots$）处有 $n_j\geq0$ 个声子，则
\begin{equation*}
E=\sum_j n_j\,v|j|k_0,\qquad K=\sum_j n_j\,j\,k_0 .
\end{equation*}
要求 $E=nk_0v$、$K=(n-2m)k_0$ 即
\begin{equation*}
\sum_{j>0}j\,(n_j+n_{-j})=n,\qquad
\sum_{j>0}j\,(n_j-n_{-j})=n-2m
\quad\Longrightarrow\quad
\sum_{j>0}j\,n_j=n-m,\quad \sum_{j>0}j\,n_{-j}=m .
\end{equation*}
正动量声子的占据数组 $(n_1,n_2,\dots)$ 与把整数 $n-m$ 分割为若干正整数之和的方式一一对应；负动量声子同理对应 $m$ 的一个分割。注意不同 $(j,k)$ 的多重集由占据数唯一确定，故
\begin{equation*}
\#\{\text{流体动力学态}\}=p_{n-m}\cdot p_m ,
\end{equation*}
其中 $p_n$ 是分割数（$(p_0,p_1,p_2,p_3,p_4,p_5)=(1,1,2,3,5,7)$）。(a)、(b) 分别是 $m=0$ 与 $m=n$ 的特例：数目为 $p_np_0=p_n$。

\textbf{费米子一侧。}取 $N$ 个无自旋费米子，$k_F=Nk_0/2$，基态把 $\{|-k_F\rangle,|{-k_F+k_0}\rangle,\dots,|k_F-k_0\rangle\}$ 逐占据。任何低能（粒子--空穴）激发都可以描述为：把 $n$ 个费米子从费米面下方搬到上方（同时也可以在两侧各自内部搬运，这不改变下述计数，只对应同一动量--能量的不同实现；等价地，直接考虑跨费米面的搬运即可覆盖全部态，因为两侧内部的重组可由跨面搬运的组合实现，而固定粒子数下两类计数一致）。设第 $i$ 个空穴在 $k_F-a_ik_0$（$a_i\geq1$），第 $j$ 个粒子在 $k_F+b_jk_0$（$b_j\geq1$），$i,j=1,\dots,n$。则
\begin{equation*}
E=\sum_{i,j}(a_i+b_j)\,k_0v_F=\Big(\sum a_i+\sum b_j\Big)k_0v_F,\qquad
K=\Big(\sum b_j-\sum a_i\Big)k_0 .
\end{equation*}
与流体动力学侧相同的要求给出 $\sum a_i=m$、$\sum b_j=n-m$。把 $m$ 写成 $n$ 个正整数（允许重复、不计顺序，多余的取 1——多余的 1 对应被搬到恰在费米面上一格的粒子，是允许的简并）之和的方式数，等于把 $m$ 分割为至多 $n$ 部分的方式数 $p_{\leq n}(m)$；由于 $n\geq m$，$p_{\leq n}(m)=p_m$。同理 $\sum b_j$ 给出 $p_{n-m}$。所以
\begin{equation*}
\boxed{\;\#\{\text{能量 } nk_0v_F,\ \text{动量 }(n-2m)k_0\text{ 的费米子粒子--空穴态}\}=p_{n-m}\,p_m\;}
\end{equation*}
与流体动力学模型完全相同。两侧的态逐一对应，能量与动量也逐一相等，故一维流体动力学模型（等价地，其 XY/玻色化描述）在小动量、低能下忠实地重现了费米子系统的全部粒子--空穴激发谱。

\emph{（验证：$n=5,m=0$ 时费米子侧五个跨面搬运对应能量 $5k_0v_F$、动量 $5k_0$ 的组态数为 $p_5=7$：$5;\,4{+}1;\,3{+}2;\,3{+}1{+}1;\,2{+}2{+}1;\,2{+}1{+}1{+}1;\,1{+}\dots{+}1$。）}

% ---------------------------------------------------------------
\section*{问题 5.1.4}
\begin{problem}
(a) 把自由费米子的压缩率 $\chi(\omega,\pmb{k})=-\Pi^{00}(\omega,\pmb{k})$（式 (4.3.2)）代入流体动力学拉格朗日量 (5.1.2) 的运动方程，求 $d=1,2,3$ 维密度模式的衰减率并与频率比较，判断是否为良好定义的声子。(b) 改为：由流体动力学模型自身计算密度响应，令其拟合式 (4.3.2)，反解 $\chi(\omega,\pmb{k})$；(a) 的结论是否改变？
\end{problem}
\solution

\textbf{(a) 朴素处理。}由 (5.1.3)，变分得运动方程（对 $\rho_{\pmb{k}}$）：
\begin{equation*}
\frac{m}{\rho_0\pmb{k}^{2}}\,\ddot{\rho}_{\pmb{k}}+\frac{1}{\chi}\rho_{\pmb{k}}=0
\quad\Longrightarrow\quad
\omega^{2}=\frac{\rho_0}{m\chi}\,\pmb{k}^{2}.
\end{equation*}
若 $\chi$ 依赖频率且 $\chi=\chi'+\mathrm{i}\chi''$（$\chi''\ll\chi'$），令 $\omega=v\pmb{k}(1+\mathrm{i}\gamma)$、$v^{2}=\rho_0/m\chi'$，代入得 $\omega^{2}\chi(\omega,\pmb{k})=\rho_0\pmb{k}^{2}/m$ 的一阶展开 $\gamma=\chi''(\omega,\pmb{k})/2\chi'$，即衰减率与频率比为
\begin{equation*}
\Gamma=-\operatorname{Im}\omega=\frac{\omega\,\chi''}{2\chi'} .
\end{equation*}
由式 (4.3.2)，$\Pi^{00}$ 的虚部集中在粒子--空穴连续谱 $|\omega|<v_Fk$ 内，且 $T\to0$、$\omega\ll v_Fk$ 时

$d=3$：$\chi'=k_Fm/2\pi^{2}$，$\chi''=\dfrac{\omega}{v_Fk}\dfrac{k_F^{2}}{4\pi v_F}$，故
\begin{equation*}
\Gamma=\frac{\omega\chi''}{2\chi'}=\frac{\pi\,\omega^{2}}{4v_Fk}
\quad\Longrightarrow\quad
\frac{\Gamma}{\omega}=\frac{\pi}{4\sqrt3}\approx0.45\quad(\omega=vk,\ v=v_F/\sqrt3);
\end{equation*}
$d=2$：$\chi'=m/2\pi$，$\chi''=\dfrac{m\omega}{2\pi v_Fk}$，故
\begin{equation*}
\Gamma=\frac{\omega^{2}}{2v_Fk}
\quad\Longrightarrow\quad
\frac{\Gamma}{\omega}=\frac{1}{2\sqrt2}\approx0.35\quad(v=v_F/\sqrt2);
\end{equation*}
$d=1$：$\Pi^{00}=\dfrac{v_F\pmb{k}^{2}}{\pi[(\omega+\mathrm{i}0^{+})^{2}-v_F^{2}\pmb{k}^{2}]}$ 的虚部全是 $\delta(\omega\mp v_Fk)$，即全部谱权重恰好集中在 $\omega=\pm v_F\pmb{k}$——而 $d=1$ 时模式频率 $v\pmb{k}=v_F\pmb{k}$ 正好落在连续谱边上，微扰公式 $\Gamma=\omega\chi''/2\chi'$ 在此处发散，不适用。

结论：$d=2,3$ 时衰减率与频率同量级（$\Gamma\sim0.4\,\omega$），密度模式\textbf{不是}良好定义的声子；$d=1$ 时朴素处理在奇异点失效，须如 (b) 自洽处理。

\textbf{(b) 自洽处理。}对外势 $V$，流体动力学模型 (5.1.2)（加 $\int\rho V$ 耦合）的运动方程为 $\frac{m}{\rho_0\pmb{k}^{2}}\ddot{\rho}_{\pmb{k}}+\frac{1}{\chi}\rho_{\pmb{k}}=V_{\pmb{k}}$，故模型自身的密度响应为
\begin{equation*}
\rho_{\pmb{k}}=\frac{V_{\pmb{k}}}{\chi^{-1}-m\omega^{2}/\rho_0\pmb{k}^{2}}
\equiv \chi_{\mathrm{hydro}}(\omega,\pmb{k})\,V_{\pmb{k}} .
\end{equation*}
令它拟合费米子结果 $\rho_{\pmb{k}}=-\Pi^{00}(\omega,\pmb{k})V_{\pmb{k}}$，反解出模型参数应取
\begin{equation*}
\frac{1}{\chi(\omega,\pmb{k})}=\frac{m\omega^{2}}{\rho_0\pmb{k}^{2}}-\frac{1}{\Pi^{00}(\omega,\pmb{k})} .
\end{equation*}
$d=1$：$\Pi^{00}=\frac{v_Fk^{2}}{\pi[(\omega+\mathrm{i}0)^{2}-v_F^{2}k^{2}]}$，且 $\rho_0=k_F/\pi$、$m/\rho_0=\pi/v_F$，故
\begin{equation*}
\frac{1}{\chi(\omega,\pmb{k})}
=\frac{\pi\omega^{2}}{v_Fk^{2}}-\frac{\pi[(\omega+\mathrm{i}0)^{2}-v_F^{2}k^{2}]}{v_Fk^{2}}
=\pi v_F\quad(\text{至无穷小虚部}),
\end{equation*}
即 $\chi(\omega,\pmb{k})=1/\pi v_F=N(E_F)$ 为\textbf{实常数}：$\Pi^{00}$ 的全部频率依赖恰好被拉格朗日量的惯性项 $m\omega^{2}/\rho_0k^{2}$ 抵消。于是 $d=1$ 模式 $\omega=v_Fk$ 无衰减——\textbf{(a) 中关于 $d=1$ 的发散结论不再成立}：一维密度模式是严格良好定义的声子（这正是玻色化在一维成立的原因：单一玻色模式穷尽了 $k\approx0$ 的全部谱权重）。

$d=2,3$：长波极限下 $m\omega^{2}/\rho_0k^{2}\propto(\omega/v_Fk)^{2}(k/k_F)\ll1/\Pi^{00}$，惯性项可忽略，
\begin{equation*}
\chi(\omega,\pmb{k})\simeq-\Pi^{00}(\omega,\pmb{k}),
\end{equation*}
与 (a) 所用假设一致。故 \textbf{(a) 在 $d>1$ 的结论不变}：衰减率 $\Gamma/\omega\sim0.4$，密度模式被粒子--空穴连续谱强烈阻尼，不是良好定义的声子；流体动力学描述在 $d>1$ 只描述平均的密度涨落。

% ---------------------------------------------------------------
\section*{问题 5.1.5}
\begin{problem}
流体动力学哈密顿量 $H=\int\mathrm{d}^dx\,[m\pmb{j}^2/2\rho_0+\rho^2/2\chi]$，设 $\partial_{\pmb{x}}\times\pmb{j}=0$，以 $\pmb{j}=\partial_{\pmb{x}}\phi$。(1) 以 $(\rho,C\phi)$ 为正则坐标--动量，求使流守恒成为哈密顿方程的 $C$；(2) 证明该正则对重现拉格朗日量 (5.1.2)；(3) 以 $(C\phi,\rho)$ 为正则对，证明重现 XY 模型拉格朗日量；(4) 由 $\pmb{k}=0$ 模的粒子数量子化求 $\phi$ 的周期条件；(5) 比较两模型算出的密度关联函数。
\end{problem}
\solution

\textbf{(1) 求 $C$。}以 $q_{\pmb{x}}=\rho(\pmb{x})$，$p_{\pmb{x}}=C\phi(\pmb{x})$，则
\begin{equation*}
\frac{\delta H}{\delta p_{\pmb{x}}}=\frac1C\frac{\delta H}{\delta\phi(\pmb{x})}
=\frac1C\Big(-\frac{m}{\rho_0}\partial_{\pmb{x}}^{2}\phi\Big)
=-\frac{m}{C\rho_0}\,\partial_{\pmb{x}}\pmb{j}.
\end{equation*}
哈密顿方程 $\dot q=\delta H/\delta p$ 给出 $\dot{\rho}=-\frac{m}{C\rho_0}\partial_{\pmb{x}}\pmb{j}$。要与流守恒 $\dot{\rho}=-\partial_{\pmb{x}}\pmb{j}$ 一致，须
\begin{equation*}
\boxed{\;C=\frac{m}{\rho_0}\;}
\end{equation*}

\textbf{(2) 重现 (5.1.2)。}第二个哈密顿方程给 $C\dot\phi=\dot p=-\delta H/\delta\rho=-\rho/\chi$，即 $\dot\phi=-\rho_0\rho/m\chi$。消去 $\phi$：$\ddot\rho=-\partial_{\pmb{x}}\dot{\pmb{j}}=-\partial_{\pmb{x}}\partial_x\dot\phi=\frac{\rho_0}{m\chi}\partial_x^{2}\rho$，即波动方程
\begin{equation*}
\ddot{\rho}-v^{2}\partial_{\pmb{x}}^{2}\rho=0,\qquad v^{2}=\frac{\rho_0}{m\chi},
\end{equation*}
这正是 (5.1.3) 的运动方程（对 (5.1.3) 变分得 $\frac{m}{\rho_0}(-\partial^{2})^{-1}\ddot\rho+\frac{\rho}{\chi}=0$，作用 $-\partial^{2}$ 后相同）。

再验证勒让德变换本身：$L=p\dot q-H=\int\mathrm{d}^dx\,[C\phi\dot\rho-m(\partial\phi)^{2}/2\rho_0-\rho^{2}/2\chi]$。由 $\dot\rho=-\frac{C\rho_0}{m}\partial_x^{2}\phi$ 解出 $\phi=-\frac{C\rho_0}{m}(-\partial^{2})^{-1}\dot\rho$，代入并分部积分，利用 $C=m/\rho_0$：
\begin{equation*}
L=\frac{C\rho_0}{2m}\int\mathrm{d}^dx\,\dot\rho(-\partial_{\pmb{x}}^{2})^{-1}\dot\rho-\int\mathrm{d}^dx\,\frac{\rho^{2}}{2\chi}
=\int\mathrm{d}^dx\,\Big(\frac{m}{2\rho_0}\dot\rho\frac{1}{-\partial_{\pmb{x}}^{2}}\dot\rho-\frac{\rho^{2}}{2\chi}\Big),
\end{equation*}
即拉格朗日量 (5.1.3)（等价于 (5.1.2)）。

\textbf{(3) 交换正则对。}取 $q=C\phi$，$p=\rho$。哈密顿方程：
\begin{equation*}
C\dot\phi=\frac{\delta H}{\delta\rho}=\frac{\rho}{\chi},\qquad
\dot\rho=-\frac1C\frac{\delta H}{\delta\phi}=\frac{m}{C\rho_0}\partial_x^{2}\phi=\frac{m}{C\rho_0}\partial_x\pmb{j},
\end{equation*}
二者联立仍给出 $\ddot\phi-v^{2}\partial_x^{2}\phi=0$。勒让德变换：$L=C\phi\dot\rho-H$，由 $\rho=C\chi\dot\phi$ 消去 $\rho$：
\begin{equation*}
L=\int\mathrm{d}^dx\,\Big[\frac{C^{2}\chi}{2}\dot\phi^{2}-\frac{m}{2\rho_0}(\partial_{\pmb{x}}\phi)^{2}\Big].
\end{equation*}
定义 $\theta=D\phi$ 并要求 $\theta$ 项系数为 $\chi/2$：由 $\frac{C^{2}\chi}{2D^{2}}=\frac{\chi}{2}$ 得 $D=C=m/\rho_0$。于是
\begin{equation*}
L=\int\mathrm{d}^dx\,\frac{\chi}{2}\Big[(\partial_t\theta)^{2}-v^{2}(\partial_{\pmb{x}}\theta)^{2}\Big],
\qquad v^{2}=\frac{m/\rho_0}{C^{2}\chi}\cdot\frac{m}{1}\cdot\frac{1}{1}=\frac{\rho_0}{m\chi},
\end{equation*}
即标准 XY 模型拉格朗日量，且声速 $v=\sqrt{\rho_0/m\chi}$ 与式 (5.1.5) 一致。故取 $\theta=(m/\rho_0)\phi$ 后，流体动力学模型就是 XY 模型。

\textbf{(4) $\phi$ 的周期性。}考虑 $\pmb{k}=0$ 模。由正则对易关系 $[\rho(\pmb{x}),C\phi(\pmb{x}')]=\mathrm{i}\delta(\pmb{x}-\pmb{x}')$ 得
\begin{equation*}
[\hat N,\phi(\pmb{x})]=\Big[\!\int\mathrm{d}^dx'\,\rho(\pmb{x}'),\phi(\pmb{x})\Big]=\frac{\mathrm{i}}{C},
\qquad\Longrightarrow\quad [\hat N,C\phi]=\mathrm{i}.
\end{equation*}
$\hat N$（总粒子数）的谱是整数，其共轭变量必为周期 $2\pi$ 的角度变量。$C\phi$ 与 $\hat N$ 恰有角度--粒子数的正则对易关系，故 $C\phi\sim C\phi+2\pi$ 描述同一物理点：
\begin{equation*}
\boxed{\;\phi\sim\phi+\frac{2\pi}{C}=\frac{2\pi\rho_0}{m}\;}
\end{equation*}
经过标度 $D\phi=(m/\rho_0)\phi=\theta$，$D\phi$ 正是周期为 $2\pi$ 的 XY 角度场 $\theta$。

\textbf{(5) 密度关联的比较。}两模型在经典层面等价（同一拉格朗日量的两种正则坐标），量子化后由 (3) 的正则变换相联系，$\rho=-\chi\partial_t\theta=-\chi(m/\rho_0)\dot\phi$。因此用同一关联（例如一维编时）
\begin{equation*}
\langle\rho(\tau)\rho(0)\rangle=\chi^{2}\partial_\tau^{2}\langle\theta(\tau)\theta(0)\rangle
=\frac{\chi v}{4\pi}\frac{1}{(v\tau)^{2}}\cdot 2\ \text{型幂律}
\end{equation*}
从 (5.1.2) 与从 XY 模型算出的密度关联函数\textbf{完全相同}（均为高斯理论，$G_\theta(\tau,x)=\frac{\mathrm{i}}{4\pi v\chi}\ln\frac{l^{2}}{x^{2}+v^{2}\tau^{2}}$ 的二阶时间导数给出 $\langle\rho\rho\rangle\propto\chi v\,l^{2}/(x^{2}+v^{2}\tau^{2})^{2}$ 型结果）。评论：该等价性是精确的高斯层面等价；差别只出现在流体动力学变量表达不了的量上——例如 XY 模型中 $\mathrm{e}^{\mathrm{i}\theta}$ 型的顶点算符（涡旋/大动量激发）在流体动力学变量 $(\rho,\phi)$ 中没有局域对应，这正是 5.1.6 要用周期性（紧致性）找回的信息。

% ---------------------------------------------------------------
\section*{问题 5.1.6}
\begin{problem}
利用流体动力学模型与 XY 模型的关系（$\theta=(m/\rho_0)\phi$，$\phi$ 周期为 $2\pi\rho_0/m$），证明一维流体动力学模型可以重现一维费米子系统中 $\pm2k_F,\pm4k_F,\dots$ 附近的低能激发（见 3.6 节）。
\end{problem}
\solution

\textbf{费米子一侧的图像。}圆周长 $L$ 上动量格点间距为 $k_0=2\pi/L$，基态占据 $\{k\in[-k_F,k_F)\}$，$k_F=\rho_0\pi$，$2k_F=2\pi\rho_0$。考虑“整体平移”组态：把每个被占据动量都平移 $nk_0$：
\begin{equation*}
|W_n\rangle:\quad k_i\mapsto k_i+nk_0\quad(\forall\ \text{被占据的 }k_i).
\end{equation*}
它仍是一个 Slater 行列式，粒子数不变；其总动量为
\begin{equation*}
K_n=N\,nk_0=\frac{2\pi nN}{L}=2\pi n\rho_0=2nk_F ,
\end{equation*}
而能量代价只有质心（伽利略助推）能量
\begin{equation*}
\Delta E=\frac{K_n^{2}}{2mN}\xrightarrow{\ N\to\infty\ }0 ,
\end{equation*}
（费米海对称，$\sum_ik_i=0$，交叉项消失）。在 $|W_n\rangle$ 之上再叠加任意小动量粒子--空穴对（能量 $v_F|q|$），便得到动量 $2nk_F$ 附近的完整低能粒子--空穴连续谱。同理 $-nk_0$ 平移给出 $-2nk_F$。

\textbf{流体动力学/XY 一侧。}XY 模型是紧致的：$\theta$ 有缠绕数为 $n$ 的位形 $\theta(x)=2\pi nx/L+\text{涨落}$。由 (5.1.5) 题，$\theta=(m/\rho_0)\phi$，故缠绕位形对应 $\phi(L)-\phi(0)=2\pi n\rho_0/m$。流体动量密度为 $\pmb{j}=\partial_x\phi$，总动量
\begin{equation*}
P=m\int_0^L\mathrm{d}x\,\partial_x\phi=m\big[\phi(L)-\phi(0)\big]=2\pi n\rho_0=2nk_F ,
\end{equation*}
即 XY 模型的第 $n$ 缠绕扇区恰好携带费米子“整体平移”激发的动量 $2nk_F$；而缠绕扇区上的小涨落（声子 $\alpha_{\pmb{k}}$）给出能量 $v|\pmb{q}|$ 的激发，能量最低的缠绕组态本身能量 $\to0$（由两侧模型谱的逐一对应，即问题 5.1.3，此能量即 $K_n^{2}/2mN$）。

\textbf{自洽性检验（密度关联）。}XY 模型中制造缠绕的算符是顶点算符 $O^{n}\sim\mathrm{e}^{\mathrm{i}2n\theta}$ 型（把 $\theta$ 缠绕 $n$ 次），其关联在 $1+1$ 维按
\begin{equation*}
\langle O^{n}(x,t)O^{n\dagger}(0,0)\rangle\sim\Big(\frac{l^{2}}{|x^{2}-v^{2}t^{2}|}\Big)^{2n^{2}\pi\chi v}
\end{equation*}
衰减，其中 $\chi=N(E_F)=1/\pi v_F$、$v=v_F$，指数 $=2n^{2}$。因此密度关联含 $k\sim2nk_F$ 的分量 $\cos(2nk_Fx)/|x-vt|^{2n^{2}}$——与自由费米子密度关联（由 $\mathcal{G}(x)\mathcal{G}(-x)$ 的幂次展开）在 $\pm2nk_F$ 处 $\sim(x-\mathrm{i}vt)^{-2n^{2}}$ 的奇异性完全一致。

结论：流体动力学模型经 $\theta=(m/\rho_0)\phi$ 等同于紧致 XY 模型，其缠绕扇区（涡旋/顶点算符）精确重现一维费米子系统 $\pm2k_F,\pm4k_F,\dots$ 附近动量大、能量低的激发及相应的密度关联奇异性。

% ---------------------------------------------------------------
\section*{问题 5.1.7}
\begin{problem}
设一维芯电子--传导电子耦合为 $\int\mathrm{d}x\,V(x)\rho(x)$，$V(x)=V/|x|^{\gamma}$，$0<\gamma<1$。求芯电子格林函数 $\mathcal{G}_C(\tau)$ 的长时间行为及发射率 $A_{em}(\omega)$。
\end{problem}
\solution

连链展开（同正文）给
\begin{equation*}
\mathcal{G}_C(\tau)\mathrm{e}^{-E_C\tau}
=\exp\Big[\frac{1}{2}\int_0^{\tau}\!\!\int_0^{\tau}\mathrm{d}\tau_1\mathrm{d}\tau_2\,F(\tau_1-\tau_2)\Big],\qquad
F(s)\equiv\int\!\!\int\mathrm{d}x\mathrm{d}y\,V(x)V(y)\langle\rho(s,x-y)\rho(0,0)\rangle .
\end{equation*}
用动量空间计算 $F(s)$。自由费米子密度关联的小动量部分傅里叶变换为
\begin{equation*}
D(s,k)=\int\mathrm{d}X\,\mathrm{e}^{-\mathrm{i}kX}\frac{N^{2}(E_F)v^{2}}{4}\Big[\frac{1}{(vs+\mathrm{i}X)^{2}}+\frac{1}{(vs-\mathrm{i}X)^{2}}\Big]
=\frac{N^{2}\pi v^{2}}{2}\,|k|\,\mathrm{e}^{-|k|vs}\qquad(s>0).
\end{equation*}
（用 $\int\mathrm{d}X\,\mathrm{e}^{-ikX}/(vs+\mathrm{i}X)^{2}=2\pi k\,\mathrm{e}^{-kvs}\Theta(k)$。）长程势的傅里叶分量（$0<\gamma<1$）
\begin{equation*}
V_k=\int\mathrm{d}x\,\frac{V}{|x|^{\gamma}}\mathrm{e}^{-\mathrm{i}kx}=2V\,\Gamma(1-\gamma)\sin\frac{\pi\gamma}{2}\,|k|^{\gamma-1}
\quad\Longrightarrow\quad |V_k|^{2}\propto|k|^{2\gamma-2},
\end{equation*}
在小 $k$ 处发散——这正是长程势的新效应。于是
\begin{equation*}
F(s)=\int\frac{\mathrm{d}k}{2\pi}\,|V_k|^{2}D(s,k)
=2N^{2}V^{2}\Gamma^{2}(1-\gamma)\sin^{2}\!\frac{\pi\gamma}{2}\;\Gamma(2\gamma)\,v^{2-2\gamma}\,s^{-2\gamma}
\equiv \frac{c}{s^{2\gamma}} ,
\end{equation*}
（积分 $\int_0^\infty\mathrm{d}k\,k^{2\gamma-1}\mathrm{e}^{-kvs}=\Gamma(2γ)/(vs)^{2γ}$；$s\lesssim\tau_l\sim1/E_F$ 处需紫外截断，其贡献只给出发散的、可并入 $E_C$ 重整化的 $\tau$ 线性项）。再利用幂次恰当地分离“平庸”部分：对 $2\gamma>1$，$\int_0^\infty\mathrm{d}s\,F(s)$ 收敛，
\begin{equation*}
W(\tau)\equiv2\int_0^{\tau}\mathrm{d}s\,(\tau-s)F(s)
=2\tau\!\int_0^{\tau}\!F-2\!\int_0^{\tau}\!sF(s)\,\mathrm{d}s
=\underbrace{2\tau\!\int_0^{\infty}\!\!\mathrm{d}s\,F(s)}_{\text{线性}\to E_C\text{ 重整化}}-\frac{2c}{2-2\gamma}\,\tau^{2-2\gamma}+O(\tau^{1-2\gamma}\tau\text{ 以下}),
\end{equation*}
其中用到 $\int_\tau^\infty s\,F(s)\mathrm{d}s=\frac{c}{2-2\gamma}\tau^{2-2\gamma}$。故长时间行为（把平庸的指数因子并入重整化后的 $\tilde E_C$）为
\begin{equation*}
\boxed{\;\mathcal{G}_C(\tau)\sim\mathrm{e}^{\tilde E_C\tau}\,
\exp\!\Big[-\alpha_\gamma\,\tau^{\,2-2\gamma}\Big],\qquad
\alpha_\gamma=\frac{2N^{2}(E_F)V^{2}\Gamma^{2}(1-\gamma)\sin^{2}\frac{\pi\gamma}{2}\,\Gamma(2\gamma)\,v^{2-2\gamma}}{1-\gamma}>0\;}
\end{equation*}
即幂律 $(\tau_l/\tau)^{\alpha}$ 被\textbf{拉伸指数}取代；此外势的有限 $2k_F$ 分量还贡献一个次要的幂律因子 $(\tau_l/\tau)^{|V_{2k_F}|^{2}N^{2}(E_F)/2}$。物理上：长程势的 $|V_k|^{2}\sim k^{2\gamma-2}$ 在小 $k$ 增强了正交性灾变，交叠衰减快于任何幂律。

讨论：$2\gamma<1$（$\gamma<1/2$）时 $\int_0^\infty F\,\mathrm{d}s$ 发散，$W(\tau)\sim+\frac{2c}{(1-2\gamma)(2-2\gamma)}\tau^{2-2\gamma}$ 为发散的正项——极化费米海所需能量发散，弱势微扰论失效（长程序被长程势强烈修正）；上式仅在 $1/2<\gamma<1$ 内受控。

\textbf{发射率。}对 $\mathcal{G}_C(t)\sim\mathrm{e}^{-\alpha_\gamma|t|^{\mu}}\mathrm{e}^{\mathrm{i}\tilde E_Ct}$（$\mu=2-2\gamma<1$）作傅里叶分析：把指数展开 $\mathrm{e}^{-\alpha\tau^\mu}=\sum_n(-\alpha)^n\tau^{\mu n}/n!$，逐项变换给出在阈值 $\Omega\equiv E_C-\omega\to0^{+}$ 处的本质奇异性
\begin{equation*}
\boxed{\;A_{em}(\omega)\sim\Theta(E_C-\omega)\,
\exp\!\Big[-\frac{c_0}{(E_C-\omega)^{\frac{2-2\gamma}{2\gamma-1}}}\Big]\;}
\end{equation*}
（$c_0\propto\alpha_\gamma^{1/(2\gamma-1)}>0$ 由鞍点确定），即幂律边缘奇异性 $(E_C-\omega)^{\alpha}$ 被完全抹平：长程势使 X 射线吸收边比任何幂律都更强烈地压低。

% ---------------------------------------------------------------
\section*{问题 5.1.8}
\begin{problem}
芯电子经 $\delta$ 势 $H_I=V_C\rho(0)$ 与一维相互作用玻色系统耦合，玻色系统用其有效 XY 模型 (3.3.23) 描述。玻色子密度关联在 $k=K_n=2\pi n\rho$ 处含奇异性（式 (3.6.2)），标度维数 $\Delta_n=n^{2}\pi\chi v$。对给定 $(\chi,v)$ 判断哪个奇异性控制 $\mathcal{G}_C(\tau)$ 的长时间行为并求出该行为。
\end{problem}
\solution

由 $\rho=\rho_0-\chi\dot\theta$（虚时 $\rho=\rho_0+\mathrm{i}\chi\partial_\tau\theta$），
\begin{equation*}
\mathcal{G}_C(\tau)=\mathrm{e}^{-V_C\rho_0\tau}\,
\big\langle\mathrm{e}^{\mathrm{i}V_C\chi\theta(\tau,0)}\mathrm{e}^{-\mathrm{i}V_C\chi\theta(0,0)}\big\rangle
=\mathrm{e}^{-V_C\rho_0\tau}\Big(\frac{l^{2}}{v^{2}\tau^{2}}\Big)^{\frac{\chi V_C^{2}}{4\pi v}}
\end{equation*}
这是“光滑密度道”（标度维数 $\Delta_0=1$，$\langle\dot\theta\dot\theta\rangle\sim1/\tau^{2}$，给对数 $\to$ 幂律）的贡献，幂指数 $\chi V_C^{2}/2\pi v$。

式 (3.6.2) 表明密度算符还含大动量道 $\sum_nC_nO_v^{n}$（$K_n=2\pi n\rho$），其关联
\begin{equation*}
\langle O_v^{n}(\tau)O_v^{n\dagger}(0)\rangle\propto\tau^{-2\Delta_n},\qquad \Delta_n=n^{2}\pi\chi v .
\end{equation*}
各道在连链展开中相加：第 $n$ 道对指数的贡献为 $W_n=|V_CC_n|^{2}\int_0^{\tau}\mathrm{d}s\,(\tau-s)c_n s^{-2\Delta_n}$。与 5.1.7 题相同的分析给出（$\Delta_n>1/2$ 时受控）：
\begin{equation*}
W_n=\text{线性项（平庸）}
-\frac{2|V_CC_n|^{2}c_n}{2-2\Delta_n}\,\tau^{\,2-2\Delta_n}\quad(1/2<\Delta_n<1),
\qquad
W_n\supset 2|V_CC_n|^{2}c_n\ln\frac{\tau}{\tau_l}\ \text{型幂律}\quad(\Delta_n=1),
\end{equation*}
而 $\Delta_n>1$ 时该道在长时间只贡献可并入 $E_C$ 的平庸项，无奇异性。长时间行为由\textbf{标度维数最小的道}控制，即 $n=\pm1$：$\Delta_1=\pi\chi v$。于是按 $(\chi,v)$ 分情形：

\textbf{(i) $\chi v>1/\pi$（$\Delta_1>1$）：}所有 $K_n$ 奇异性均无关（irrelevant），光滑道主导：
\begin{equation*}
\boxed{\;\mathcal{G}_C(\tau)\sim\mathrm{e}^{-V_C\rho_0\tau}\,\tau^{-\chi V_C^{2}/2\pi v}\;}
\end{equation*}
普通幂律正交性灾变。

\textbf{(ii) $\pi\chi v=1$（边际）：}$n=1$ 道给对数，幂律指数被加强：
\begin{equation*}
\mathcal{G}_C(\tau)\sim\mathrm{e}^{-V_C\rho_0\tau}\,\tau^{-\left(\frac{\chi V_C^{2}}{2\pi v}+2|V_CC_1|^{2}c_1\right)} .
\end{equation*}

\textbf{(iii) $1/2<\pi\chi v<1$（$n=1$ 道相关但收敛）：}$n=1$ 涡旋道主导，
\begin{equation*}
\boxed{\;\mathcal{G}_C(\tau)\sim\exp\!\Big[-\frac{2|V_CC_1|^{2}c_1}{2-2\pi\chi v}\,\tau^{\,2-2\pi\chi v}\Big]\;}
\end{equation*}
拉伸指数衰减，强于任何幂律。

\textbf{(iv) $\pi\chi v<1/2$：}$\int_0 s\,s^{-2\Delta_1}\mathrm{d}s$ 在红外发散，连链展开指数增长——弱耦合微扰论失效（涡旋算符高度相关，芯势引起强重构）。

物理总结：玻色系统中 $K_n$ 奇异性的标度维数 $\Delta_n=n^{2}\pi\chi v$ 随 $\chi v$ 减小而减小；$\chi v>1/\pi$ 时大动量道无关、芯电子行为如自由颗粒情形只受光滑密度涨落影响；$\chi v\leq1/\pi$ 时 $K_1$ 奇异性控制甚至摧毁幂律边缘奇异性。

% ---------------------------------------------------------------
\section*{问题 5.2.1}
\begin{problem}
一维格点电子系统 $H=-t\sum_i(c^{\dagger}_{\alpha,i}c_{\alpha,i+1}+h.c.)+V\sum_i(\sum_\alpha c^{\dagger}_{\alpha,i}c_{\alpha,i})^{2}$，$N<N_{\mathrm{site}}$。(1) 求 $V=0$ 的 $\epsilon_k$；(2) 用 Hartree--Fock 求 $E(N_{\uparrow},N_{\downarrow})$；(3) 在 $N_{\uparrow}-N_{\downarrow}=0$ 邻域判断稳定性并求 $V_0$；(4) 求自旋磁化率 $\chi$ 及其在 $V\to V_0$ 的行为；(5) 固定 $N$ 极小化，求铁磁临界 $V_c$ 与基态总自旋 $S_{z0}$；(6) 对 $V\lessgtr V_c$ 求 $\Sigma_{\pmb{k}\alpha}$ 与激发谱；(7) 求翻转一个自旋的最小能量并评论。
\end{problem}
\solution

\textbf{(1)} 色散关系为 $\epsilon_k=-2t\cos k$（晶格常数取 1）。

\textbf{(2)} 在位势对应 $V_q=V$。以占据数 $n_{\pmb{k}\alpha}\in\{0,1\}$ 的 Slater 行列式为尝试态，由式 (5.2.2)（$V_0\to V$）：
\begin{equation*}
E=\sum_{\pmb{k},\alpha}n_{\pmb{k}\alpha}\,\epsilon_k
+\frac{V}{2N_{\mathrm{site}}}N^{2}
-\frac{V}{2N_{\mathrm{site}}}\big(N_{\uparrow}^{2}+N_{\downarrow}^{2}\big),
\qquad N_\alpha=\sum_{\pmb{k}}n_{\pmb{k}\alpha}.
\end{equation*}
（Hartree 项 $\frac{V}{2N_{\mathrm{site}}}N^{2}$ 与总密度无关地极化；Fock 项只作用在同自旋对之间。）动能显式：自旋 $\alpha$ 填充 $|k|\leq k_{F\alpha}=\pi N_\alpha/N_{\mathrm{site}}$，
\begin{equation*}
E_{\mathrm{kin},\alpha}=\frac{N_{\mathrm{site}}}{2\pi}\int_{-k_{F\alpha}}^{k_{F\alpha}}(-2t\cos k)\,\mathrm{d}k
=-\frac{2N_{\mathrm{site}}t}{\pi}\sin k_{F\alpha},
\end{equation*}
\begin{equation*}
\boxed{\;E(N_{\uparrow},N_{\downarrow})=-\frac{2N_{\mathrm{site}}t}{\pi}\big[\sin(\pi N_{\uparrow}/N_{\mathrm{site}})+\sin(\pi N_{\downarrow}/N_{\mathrm{site}})\big]
+\frac{V}{2N_{\mathrm{site}}}N^{2}-\frac{V}{2N_{\mathrm{site}}}(N_{\uparrow}^{2}+N_{\downarrow}^{2})\;}
\end{equation*}

\textbf{(3)} 令 $N_{\uparrow}=N/2+s$，$N_{\downarrow}=N/2-s$（$s=(N_{\uparrow}-N_{\downarrow})/2$）。$\sin$ 展开到 $s^{2}$：$\sin(k_F{+}\delta)+\sin(k_F{-}\delta)=2\sin k_F\cos\delta\approx2\sin k_F\,[1-(\pi s/N_{\mathrm{site}})^{2}/2]$，$k_F\equiv\pi N/2N_{\mathrm{site}}$；而 $N_{\uparrow}^{2}+N_{\downarrow}^{2}=N^{2}/2+2s^{2}$。于是（动能贡献 $+\frac{2\pi t\sin k_F}{N_{\mathrm{site}}}s^{2}$，Fock 贡献 $-\frac{V}{N_{\mathrm{site}}}s^{2}$）
\begin{equation*}
E(s)=E(0)+\frac{s^{2}}{N_{\mathrm{site}}}\big(V_0-V\big)
=E(0)+\frac{(N_{\uparrow}-N_{\downarrow})^{2}}{4N_{\mathrm{site}}}\big(V_0-V\big),
\qquad
\boxed{\;V_0=2\pi t\sin k_F=2\pi t\sin\frac{\pi N}{2N_{\mathrm{site}}}\;}
\end{equation*}
$V<V_0$ 时 $s=0$（即 $N_{\uparrow}=N_{\downarrow}$）为局部极小（顺磁态稳定）；$V>V_0$ 时为局部极大。物理：Fock（同自旋吸引）能增益超过动能代价时发生铁磁失稳。

\textbf{(4)} 加磁场项 $-h(N_{\uparrow}-N_{\downarrow})=-2hs$，极小化 $E(s)-2hs$：$2s(V_0-V)/N_{\mathrm{site}}=2h$，得 $s=hN_{\mathrm{site}}/(V_0-V)$，即
\begin{equation*}
\chi=\frac{\partial (N_{\uparrow}-N_{\downarrow})}{\partial h}\Big|_{h\to0}
=\frac{2N_{\mathrm{site}}}{V_0-V}
\quad\xrightarrow{\ V\to V_0^{-}\ }\ \infty ,
\end{equation*}
磁化率发散（Weiss 增强 $\propto1/(1-V/V_0)$）；$V=0$ 时回到 Pauli 值 $2N_{\mathrm{site}}\mu_B^{2}N(E_F)$（每自旋 $N(E_F)=1/2\pi t\sin k_F$，与 $2N_{\mathrm{site}}/V_0$ 一致）。

\textbf{(5)} 固定 $N$ 全域极小化。$E(s)=-\frac{4N_{\mathrm{site}}t}{\pi}\sin k_F\cos(\pi s/N_{\mathrm{site}})-\frac{V}{N_{\mathrm{site}}}s^{2}+\text{常数}$（$0\leq s\leq N/2$）。$\mathrm{d}E/\mathrm{d}s=4t\sin k_F\sin(\pi s/N_{\mathrm{site}})-4Vs/N_{\mathrm{site}}$。展开 $E(s)=E(0)+a s^{2}+b s^{4}+\dots$，$a=(V_0-V)/N_{\mathrm{site}}$，而 $b=-\frac{\pi^{3}t\sin k_F}{6N_{\mathrm{site}}^{3}}<0$：四次项为负，故相变为\textbf{一级}——顺磁态一旦失稳即直接跳到完全极化 $s=N/2$（可验证对 $V>V_0/2$ 时 $\mathrm{d}E/\mathrm{d}s=0$ 无内部驻点）。比较完全极化能
\begin{equation*}
E_{\mathrm{fp}}=-\frac{2N_{\mathrm{site}}t}{\pi}\sin(2k_F)+\frac{V}{2}N
\end{equation*}
（Hartree 与 Fock 相消：$V N^{2}/2N_{\mathrm{site}}-VN^{2}/2N_{\mathrm{site}}=0$）与 $E(0)=-\frac{4N_{\mathrm{site}}t}{\pi}\sin k_F+\frac{V N^{2}}{4N_{\mathrm{site}}}+\frac{V N}{2}$，得
\begin{equation*}
\boxed{\;V_c=\frac{8N_{\mathrm{site}}^{2}t}{\pi N^{2}}\big[2\sin k_F-\sin 2k_F\big]
=\frac{8t}{\pi n^{2}}\Big[2\sin\frac{\pi n}{2}-\sin(\pi n)\Big],\quad n\equiv\frac{N}{N_{\mathrm{site}}}\;}
\end{equation*}
低填充 $n\to0$：$V_c\to\pi^{2}tn=V_0$，过渡趋于连续；基态为完全极化铁磁体，\textbf{$S_{z0}=N/2$}（沿 $z$ 全部同向）。

\textbf{(6)} 由式 (5.2.4)（$V_q=V$）：$\Sigma_{\pmb{k}\alpha}=-\frac{V}{N_{\mathrm{site}}}\sum_{\pmb{k}'}n_{\pmb{k}'\alpha}=-\frac{VN_{\alpha}}{N_{\mathrm{site}}}$，与 $\pmb{k}$ 无关。

$V<V_c$（顺磁，$N_\uparrow=N_\downarrow=N/2$）：$\Sigma_{\pmb{k}\alpha}=-\frac{VN}{2N_{\mathrm{site}}}$ 自旋无关，$\xi^{*}_{\pmb{k}\alpha}=-2t\cos k-\mu'-\frac{VN}{2N_{\mathrm{site}}}$：就是自由费米子谱平移一个常数；激发为费米面上下的粒子--空穴对，能量 $\epsilon_k-\epsilon_{k'}$（$k>k_F$，$k'<k_F$），无能隙。

$V>V_c$（铁磁，$N_\uparrow=N$，$N_\downarrow=0$）：$\Sigma_{\pmb{k}\uparrow}=-\frac{VN}{N_{\mathrm{site}}}$，$\Sigma_{\pmb{k}\downarrow}=0$。激发谱分两支：(i) 多数（$\uparrow$）带内的粒子--空穴激发，能量 $\epsilon_k-\epsilon_{k'}$，无能隙；(ii) 产生一个 $\downarrow$ 准粒子：能量 $\xi^{*}_{k\downarrow}-\xi^{*}_{k\uparrow}=\frac{VN}{N_{\mathrm{site}}}$，\textbf{有限能隙} $\Delta=\frac{VN}{N_{\mathrm{site}}}$。

\textbf{(7)} 翻转一个自旋（$\Delta S_z=1$）的最小能量：顺磁态中把 $\uparrow$ 费米面上的电子翻到同动量的 $\downarrow$ 空位：代价 $\Sigma_\downarrow-\Sigma_\uparrow=0$，即零能量（自旋简并，正确）。铁磁态：$\Delta=\frac{VN}{N_{\mathrm{site}}}>0$，有限。\textbf{评论：}铁磁态自发破缺自旋转动对称性，按 Goldstone 定理必有无能隙自旋波。真实铁磁体中翻一个自旋所形成的局域磁矩可以弥散到许多格点上（多自旋翻转的相干叠加＝长波长自旋波），能量 $\to0$；Hartree--Fock 把每次自旋翻转当独立粒子处理，漏掉了这种集体效应，给出的有限自旋翻转能隙 $\frac{VN}{N_{\mathrm{site}}}$ \textbf{不可信}——正确的铁磁体自旋翻转能隙为零。

% ---------------------------------------------------------------
\section*{问题 5.2.2}
\begin{problem}
用平均场近似（把 $c^{\dagger}_{\alpha\pmb{k}}c_{\beta\pmb{k}'}$ 换成其平均值）导出式 (5.2.7) 的有效哈密顿量；不同换法给出若干项贡献；补上常数项使 $\langle H_{\mathrm{eff}}\rangle=\langle H\rangle$。
\end{problem}
\solution

把相互作用写成 $H_I=\frac{1}{2N_{\mathrm{site}}}\sum_{\pmb{k}\pmb{k}'\pmb{q}}V_{\pmb{q}}\,c^{\dagger}_{\alpha\pmb{k}}c_{\alpha,\pmb{k}-\pmb{q}}\,c^{\dagger}_{\beta\pmb{k}'}c_{\beta,\pmb{k}'+\pmb{q}}$。平均场近似有两条（对易的）换法：

\textbf{直接（Hartree）道}：把每个 $c^{\dagger}c$ 换成 $\langle c^{\dagger}c\rangle+(\text{涨落})$ 并保留线性项，两个双线性算符各换一次，扣除重复计数：
\begin{equation*}
H_{\mathrm{dir}}=\frac{1}{N_{\mathrm{site}}}\sum_{\pmb{k}\pmb{k}'\pmb{q}}V_{\pmb{q}}\,
\langle c^{\dagger}_{\alpha\pmb{k}}c_{\alpha,\pmb{k}-\pmb{q}}\rangle\,c^{\dagger}_{\beta\pmb{k}'}c_{\beta,\pmb{k}'+\pmb{q}}
-\frac{1}{2N_{\mathrm{site}}}\sum V_{\pmb{q}}\langle c^{\dagger}c\rangle\langle c^{\dagger}c\rangle .
\end{equation*}
平移不变下 $\langle c^{\dagger}_{\alpha\pmb{k}}c_{\alpha,\pmb{k}-\pmb{q}}\rangle=n_{\pmb{k}\alpha}\delta_{\pmb{q}0}$，故
\begin{equation*}
H_{\mathrm{dir}}=\sum_{\pmb{k}\alpha}\big(V_0\rho_0\big)c^{\dagger}_{\alpha\pmb{k}}c_{\alpha\pmb{k}}-\text{常数}.
\end{equation*}

\textbf{交换（Fock）道}：先利用 $c_{\alpha,\pmb{k}-\pmb{q}}c^{\dagger}_{\beta\pmb{k}'}=-c^{\dagger}_{\beta\pmb{k}'}c_{\alpha,\pmb{k}-\pmb{q}}+\delta_{\alpha\beta}\delta_{\pmb{k}-\pmb{q},\pmb{k}'}$ 改写，再把 $c^{\dagger}_{\alpha\pmb{k}}c_{\beta,\pmb{k}'+\pmb{q}}$（跨顶点配对）换成其平均 $\langle c^{\dagger}_{\alpha\pmb{k}}c_{\beta,\pmb{k}'+\pmb{q}}\rangle=\delta_{\alpha\beta}n_{\pmb{k}\alpha}\delta_{\pmb{k},\pmb{k}'+\pmb{q}}$，两个交换位各换一次并扣除重复：
\begin{equation*}
H_{\mathrm{ex}}=\frac{1}{N_{\mathrm{site}}}\sum_{\pmb{k}\pmb{q}\alpha}V_{\pmb{q}}\,n_{\pmb{k}\alpha}\,
c^{\dagger}_{\alpha,\pmb{k}-\pmb{q}}c_{\alpha,\pmb{k}-\pmb{q}}
-\frac{1}{2N_{\mathrm{site}}}\sum_{\pmb{k}\pmb{q}\alpha}V_{\pmb{q}}\,n_{\pmb{k}\alpha}n_{\pmb{k}-\pmb{q},\alpha}
=\sum_{\pmb{k}\alpha}\Sigma^{\mathrm{Fock}}_{\pmb{k}\alpha}\,c^{\dagger}_{\alpha\pmb{k}}c_{\alpha\pmb{k}}-\text{常数},
\end{equation*}
其中 $\Sigma^{\mathrm{Fock}}_{\pmb{k}\alpha}=-\frac{1}{N_{\mathrm{site}}}\sum_{\pmb{q}}V_{\pmb{q}}\,n_{\pmb{k}-\pmb{q},\alpha}$，正是式 (5.2.4)。

合并（并把常数写成显式）：
\begin{equation*}
\boxed{\;H_{\mathrm{eff}}=\sum_{\pmb{k},\alpha}\big(\epsilon_{\pmb{k}}-\mu+V_0\rho_0+\Sigma^{\mathrm{Fock}}_{\pmb{k}\alpha}\big)c^{\dagger}_{\alpha\pmb{k}}c_{\alpha\pmb{k}}
+\frac{1}{2N_{\mathrm{site}}}\sum_{\pmb{k}\pmb{q}\alpha}V_{\pmb{q}}\,n_{\pmb{k}\alpha}n_{\pmb{k}-\pmb{q},\alpha}\;}
\end{equation*}
（这就是式 (5.2.7) 加常数项，因 $\rho_0V_0$ 并入 $\mu'$ 后与 (5.2.7) 一致。）

\textbf{常数项的确定。}取平均：$\langle H_{\mathrm{eff}}\rangle_{\mathrm{int}}=V_0\rho_0N-\frac{1}{N_{\mathrm{site}}}\sum V_{\pmb{q}}nn-\frac{1}{2N_{\mathrm{site}}}\sum V_{\pmb{q}}nn$。另一方面 Wick 定理给出精确的一阶期望值
\begin{equation*}
\langle H_I\rangle=\frac{V_0}{2}\rho_0N+\frac{1}{2N_{\mathrm{site}}}\sum_{\pmb{k}\pmb{q}\alpha}V_{\pmb{q}}n_{\pmb{k}\alpha}\big(1-n_{\pmb{k}-\pmb{q},\alpha}\big)
=V_0\rho_0N-\frac{1}{2N_{\mathrm{site}}}\sum_{\pmb{k}\pmb{q}\alpha}V_{\pmb{q}}n_{\pmb{k}\alpha}n_{\pmb{k}-\pmb{q},\alpha},
\end{equation*}
（用 $\sum_{\pmb{q}}V_{\pmb{q}}=N_{\mathrm{site}}V_0$）。两者比较得常数项必须取 $+\frac{1}{2N_{\mathrm{site}}}\sum_{\pmb{k}\pmb{q}\alpha}V_{\pmb{q}}n_{\pmb{k}\alpha}n_{\pmb{k}-\pmb{q},\alpha}$，于是 $\langle H_{\mathrm{eff}}\rangle=\langle H\rangle$（对平均场所用的 $n_{\pmb{k}\alpha}$ 精确成立）。注意 $\langle H_{\mathrm{eff}}\rangle$ 含常数且相互作用只计一次，而 naive 替换会把相互作用计两次——这正是需要补偿常数项的原因。

% ---------------------------------------------------------------
\section*{问题 5.3.1}
\begin{problem}
计算 Landau 费米液体的压缩率 $\chi$ 与自旋磁化率 $\chi_s$，并求 $\chi$、$\chi_s$ 与 $C_V$ 之间的比值，与自由电子系统比较。
\end{problem}
\solution

设各向同性费米液体，费米液体函数在费米面上的角平均记为
\begin{equation*}
\bar f=\int\frac{\mathrm{d}\Omega_{\pmb{k}'}}{4\pi}\,f(\pmb{k},\alpha;\pmb{k}',\beta)
\quad(\text{与 }\pmb{k},\alpha,\beta\text{ 无关})，
\end{equation*}
每自旋态密度 $N^{*}(E_F)=k_F^{2}/(2\pi^{2}v_F^{*})$（三维，$\hbar=1$），总态密度 $\nu^{*}=2N^{*}$。

\textbf{压缩率。}均匀地改变密度：$k_F\to k_F+\mathrm{d}k$。准粒子能量变化 $\mathrm{d}\xi^{*}=v_F^{*}\mathrm{d}k$；同时被移动的 $\delta n$ 壳层通过 $f$ 反作用于准粒子，$\epsilon^{T}=\xi^{*}+\frac{1}{\mathcal V}\sum f\,\delta n$。平衡时 $\epsilon^{T}(k_F)=\mu$，故
\begin{equation*}
\mathrm{d}\mu=v_F^{*}\mathrm{d}k+\frac{1}{\mathcal V}\sum_{\pmb{k}'\beta}f\,\delta n_{\pmb{k}'\beta}
=v_F^{*}\mathrm{d}k+\bar f\cdot\frac{2k_F^{2}\mathrm{d}k}{2\pi^{2}}
=v_F^{*}\mathrm{d}k+\nu^{*}\bar f\,v_F^{*}\,\mathrm{d}k ,
\end{equation*}
其中用了 $\frac{1}{\mathcal V}\sum_{\pmb{k}'\beta}\delta n_{\pmb{k}'\beta}=\frac{2k_F^{2}\mathrm{d}k}{2\pi^{2}}=\nu^{*}v_F^{*}\mathrm{d}k$。另一方面密度改变为 $\mathrm{d}n=\nu^{*}v_F^{*}\mathrm{d}k$，于是
\begin{equation*}
\chi\equiv\frac{\partial n}{\partial\mu}
=\frac{\nu^{*}v_F^{*}\,\mathrm{d}k}{v_F^{*}(1+\nu^{*}\bar f)\,\mathrm{d}k}
=\boxed{\;\frac{\nu^{*}}{1+\nu^{*}\bar f}
=\frac{2N^{*}(E_F)}{1+2N^{*}(E_F)\bar f}\;}
\end{equation*}

\textbf{自旋磁化率。}加磁场使 $\uparrow/\downarrow$ 能量移动 $\mp h$。设壳层占据移动 $\delta\langle n_{\uparrow}\rangle=+\delta$，$\delta\langle n_{\downarrow}\rangle=-\delta$。自旋转动不变下记 $f^{d}=f(\uparrow;\uparrow)$，$f^{x}=f(\uparrow;\downarrow)$，则
\begin{equation*}
\epsilon^{T}_{\uparrow}-\xi^{*}=-h+\frac{1}{\mathcal V}\sum_{\pmb{k}'}\big(f^{d}\delta+f^{x}(-\delta)\big)
=-h+N^{*}(E_F)\big(f^{d}-f^{x}\big)\delta .
\end{equation*}
平衡条件 $\epsilon^{T}_{\uparrow}=\mu$（对 $\downarrow$ 同理反号）给出 $\delta\big[1+N^{*}(f^{d}-f^{x})\big]=N^{*}h$，磁化 $M=\mu_B\,\mathrm{2}\delta$：
\begin{equation*}
\boxed{\;\chi_s=\frac{2N^{*}(E_F)\mu_B^{2}}{1+N^{*}(E_F)\big(f^{d}-f^{x}\big)}\;}
\end{equation*}
（对式 (5.3.2) 的 $f=V_0-V_{\pmb{k}-\pmb{k}'}\delta_{\alpha\beta}$：$f^{d}-f^{x}=-V_{\pmb{k}-\pmb{k}'}$，角平均 $-\bar V$；排斥作用增大 $f^{d}-f^{x}$ 不成立时即 $-\bar V<0$ 增强 $\chi_s$——与 5.5 节 RPA 定性一致。）

\textbf{与 $C_V$ 的比值。}$C_V=\frac{2\pi^{2}}{3}\cdot2N^{*}(E_F)\,T=\frac{4\pi^{2}}{3}N^{*}T$（与自由情形同式，只把 $m\to m^{*}$）。于是
\begin{equation*}
\frac{\chi}{C_V/T}=\frac{3}{2\pi^{2}}\cdot\frac{1}{1+2N^{*}\bar f},
\qquad
\frac{\chi_s}{C_V/T}=\frac{3\mu_B^{2}}{2\pi^{2}}\cdot\frac{1}{1+N^{*}(f^{d}-f^{x})},
\end{equation*}
而 Wilson 比
\begin{equation*}
R_W\equiv\frac{\chi_s/\mu_B^{2}}{\chi}
=\frac{1+2N^{*}(E_F)\bar f}{1+N^{*}(E_F)(f^{d}-f^{x})} .
\end{equation*}
\textbf{与自由电子系统比较：}$f=0$ 时 $\chi=2N^{*}(E_F)$、$\chi_s=2N^{*}(E_F)\mu_B^{2}$、$R_W=1$，且 $\chi$ 与 $C_V/T$ 只依赖同一 $m^{*}$（比值 $3/2\pi^{2}\mu_B^{2}$-型常数）；相互作用通过 $f$ 把三者\textbf{以不同方式}重整化——$C_V$ 只测 $m^{*}$，$\chi$ 测 $m^{*}$ 与电荷道 $f$ 的组合，$\chi_s$ 测 $m^{*}$ 与自旋道 $f^{d}-f^{x}$ 的组合，故 $\chi,\chi_s$ 与 $C_V$ 的比值一般偏离自由值，$R_W\neq1$。

% ---------------------------------------------------------------
\section*{问题 5.3.2}
\begin{problem}
证明式 (5.3.5) 与 (5.3.6)（数流密度及其拖曳/回流修正）。
\end{problem}
\solution

\textbf{(5.3.5)。}无外力时玻尔兹曼方程为
\begin{equation*}
\frac{\partial n}{\partial t}
+\frac{\partial n}{\partial\pmb{x}}\cdot\frac{\partial\tilde\epsilon}{\partial\pmb{k}}
-\frac{\partial n}{\partial\pmb{k}}\cdot\frac{\partial\tilde\epsilon}{\partial\pmb{x}}=I[n] .
\end{equation*}
用提示把后两项写成全散度：
\begin{equation*}
\frac{\partial n}{\partial\pmb{x}}\cdot\frac{\partial\tilde\epsilon}{\partial\pmb{k}}
-\frac{\partial n}{\partial\pmb{k}}\cdot\frac{\partial\tilde\epsilon}{\partial\pmb{x}}
=\partial_{\pmb{x}}\cdot\Big(n\frac{\partial\tilde\epsilon}{\partial\pmb{k}}\Big)
-\partial_{\pmb{k}}\cdot\Big(n\frac{\partial\tilde\epsilon}{\partial\pmb{x}}\Big).
\end{equation*}
对 $\pmb{k}$ 积分：左边第一项给 $\partial_t\int\mathrm{d}\pmb{k}\,n=\partial_t\rho(\pmb{x},t)$；第二项给 $\partial_{\pmb{x}}\cdot\pmb{J}$，其中
\begin{equation*}
\pmb{J}(\pmb{x},t)=\mathcal{V}^{-1}\sum_{\pmb{k}}n_{\pmb{k}\alpha}(\pmb{x},t)\,\frac{\partial\tilde\epsilon}{\partial\pmb{k}}
\end{equation*}
（对 $\pmb{k}$ 的全导数项因 $|\pmb{k}|\to\infty$ 处 $n\to0$ 而无边界贡献）；碰撞项因粒子数守恒 $\sum_{\pmb{k}}I[n]=0$ 而消失。得连续性方程 $\partial_t\rho+\partial_{\pmb{x}}\cdot\pmb{J}=0$，即式 (5.3.5)。

\textbf{(5.3.6)。}线性化：$n=n_0+\delta n$，$\tilde\epsilon=\xi^{*}+\frac{1}{\mathcal V}\sum f\,\delta n\equiv\xi^{*}+\delta\epsilon_{\pmb{k}}$，
\begin{equation*}
\pmb{J}=\mathcal{V}^{-1}\sum_{\pmb{k}}\delta n_{\pmb{k}}\,\partial_{\pmb{k}}\tilde\epsilon
=\mathcal{V}^{-1}\sum_{\pmb{k}}\delta n_{\pmb{k}}\,\pmb{v}^{*}(\pmb{k})
+\mathcal{V}^{-1}\sum_{\pmb{k}}n_{0,\pmb{k}}\,\partial_{\pmb{k}}\delta\epsilon_{\pmb{k}}
+O(\delta n^{2}),
\end{equation*}
（$\delta n\,\partial_k\delta\epsilon$ 为二阶，舍去）。第二项对 $\pmb{k}$ 分部积分：
\begin{equation*}
\mathcal{V}^{-1}\sum_{\pmb{k}}n_{0,\pmb{k}}\partial_{\pmb{k}}\delta\epsilon_{\pmb{k}}
=-\mathcal{V}^{-1}\sum_{\pmb{k}}\delta\epsilon_{\pmb{k}}\,\partial_{\pmb{k}}n_{0,\pmb{k}}
=+\mathcal{V}^{-1}\sum_{\pmb{k}}\delta\epsilon_{\pmb{k}}\,\pmb{v}^{*}(\pmb{k})\,\delta(\xi^{*}_{\pmb{k}}),
\end{equation*}
其中用了 $\partial_{\pmb{k}}n_{0,\pmb{k}}=\partial_{\pmb{k}}\Theta(-\xi^{*}_{\pmb{k}})=-\pmb{v}^{*}(\pmb{k})\delta(\xi^{*}_{\pmb{k}})$（提示）。代入 $\delta\epsilon_{\pmb{k}}=\frac{1}{\mathcal V}\sum_{\pmb{k}'}f(\pmb{k},\alpha;\pmb{k}',\beta)\delta n_{\pmb{k}'\beta}$ 并换哑标：
\begin{equation*}
\pmb{J}(\pmb{x},t)=\mathcal{V}^{-1}\sum_{\pmb{k}}\delta n_{\pmb{k}\alpha}\,\tilde{\pmb{v}}(\pmb{k}),
\qquad
\tilde{\pmb{v}}(\pmb{k})=\pmb{v}^{*}(\pmb{k})
+\int\frac{\mathrm{d}^{d}\pmb{k}'}{(2\pi)^{d}}\sum_{\beta}f(\pmb{k},\alpha;\pmb{k}',\beta)\,\pmb{v}^{*}(\pmb{k}')\,\delta(\xi^{*}_{\pmb{k}'}) ,
\end{equation*}
即式 (5.3.6)。物理解释：处在 $\pmb{k}$ 的准粒子通过 $f$ 使整个费米海的准粒子速度发生微小改变，费米海的这些回流也对 $\pmb{J}$ 有贡献——第二项即拖曳/回流项。

% ---------------------------------------------------------------
\section*{问题 5.3.3}
\begin{problem}
把二维转动不变系统的约化玻尔兹曼方程 (5.3.11) 与哈密顿量 (5.3.12) 推广到没有转动对称性的二维系统。
\end{problem}
\solution

无转动对称时费米面不再是圆：沿费米面参数化 $\pmb{k}_F(\theta)$（$\theta$ 沿弧长或任意单调参数），单位法向 $\hat{\pmb{k}}(\theta)$、切向 $\hat{\pmb{k}}_{\perp}(\theta)$，且
\begin{equation*}
v_F^{*}=v_F^{*}(\theta)\equiv\big|\partial_{\pmb{k}}\xi^{*}_{\pmb{k}}\big|_{\pmb{k}=\pmb{k}_F(\theta)},\qquad
k_F=k_F(\theta).
\end{equation*}
费米面位移 $h(\theta)$ 与 $\tilde\rho$ 的关系照旧：$\tilde\rho(\theta)=\frac{k_F(\theta)h(\theta)}{(2\pi)^{2}}$。把正文推导（对式 (5.3.7) 施加 $\int\frac{k\,\mathrm{d}k}{(2\pi)^{2}}$，逐条代 $\delta(|\pmb{k}|-k_F)\to\delta(|\pmb{k}|-k_F(\theta))$、$\partial_{\pmb{k}}\xi^{*}\to v_F^{*}(\theta)\hat{\pmb{k}}(\theta)$）重新执行，每一处 $k_F$、$v_F^{*}$ 变为 $\theta$ 的函数，且 $f(\theta,\theta')$ 不再只依赖 $\theta-\theta'$，线性化约化玻尔兹曼方程成为
\begin{equation*}
\boxed{\;
\frac{\partial \tilde\rho(\theta)}{\partial t}
+v_F^{*}(\theta)\,\hat{\pmb{k}}(\theta)\cdot\frac{\partial \tilde\rho(\theta)}{\partial\pmb{x}}
+\frac{k_F(\theta)}{(2\pi)^{2}}\int\mathrm{d}\theta'\,f(\theta,\theta')\,
\hat{\pmb{k}}(\theta)\cdot\frac{\partial\tilde\rho(\theta')}{\partial\pmb{x}}
=0\;}
\end{equation*}
（$\tau^{-1}\tilde\rho$ 与力项照式 (5.3.9)、(5.3.10) 同样替换）。相应地，哈密顿量 (5.3.12) 逐字推广为
\begin{equation*}
\boxed{\;
H=E_g+\int\mathrm{d}^{2}\pmb{x}\,\mathrm{d}\theta\,
\frac{2\pi^{2}v_F^{*}(\theta)}{k_F(\theta)}\,\tilde\rho^{2}(\theta)
+\frac12\int\mathrm{d}^{2}\pmb{x}\,\mathrm{d}\theta\,\mathrm{d}\theta'\,
f(\theta,\theta')\,\tilde\rho(\theta)\tilde\rho(\theta')\;}
\end{equation*}
（动能项 $\int\tilde\rho\,\frac{hv_F^{*}}{2}$ 中 $h=\frac{(2\pi)^{2}\tilde\rho}{k_F(\theta)}$。）。稳定性条件随之改为：矩阵 $M(\theta,\theta')=v_F^{*}(\theta)\delta(\theta-\theta')\cdot(2\pi)^{2}/k_F(\theta)\cdot\frac{1}{2\pi^{2}}\cdots$ 连同 $f(\theta,\theta')$ 构成的整体核正定。所有推导只用“费米面局地近似的法向位移”这一几何事实，故上述替换是系统的、不需转动对称性；集体模谱仍由（把 $v_F^{*}+\frac{k_Ff_l}{2\pi}$ 换成 $\theta$-依赖矩阵后的）$q\tilde M^{\top}K\tilde M$ 型实对称算符的本征值给出，只是 $l$ 不再是好量子数，须直接对角化积分核。

% ---------------------------------------------------------------
\section*{问题 5.3.4}
\begin{problem}
(a) 导出一维费米液体（两支费米点 $\pm k_F$）的玻色化拉格朗日量（(5.3.13) 的一维版本）。(b) 对式 (5.1.3) 以 $\rho=(2\pi)^{-1}\partial_x\varphi$ 引入 $\varphi$，求正则动量 $\tilde\varphi$、哈密顿量 $H$，并把作用量 $S=\int\mathrm{d}t(\tilde\varphi\dot\varphi-H)$ 与 (a) 比较。
\end{problem}
\solution

\textbf{(a)} 一维费米面是两点 $\pm k_F$；用 $\tilde\rho_r(x,t)$（$r=\pm$）标记两支的位移密度。仿照 (5.3.13) 的推导（对称化 $\theta\leftrightarrow r$，$\int\mathrm{d}\theta\to\sum_{r}$，$\hat{\pmb{k}}\cdot\partial_{\pmb{x}}\to r\partial_x$；取每支归一 $\tilde\rho_r=k_F\!\int\frac{\mathrm{d}k}{2\pi}\delta n_{rk}$，使 $h_r=2\pi\tilde\rho_r/k_F$）：
\begin{align*}
L&=\sum_{r}\int\mathrm{d}x\,\Big[\frac{1}{2k_F}\,r\,\partial_x\tilde\varphi_r\,\partial_t\tilde\varphi_r
-\frac{v_F^{*}}{2k_F}\,\big(r\,\partial_x\tilde\varphi_r\big)^{2}\Big]
-\frac{1}{8\pi^{2}}\int\mathrm{d}x\sum_{rr'}f_{rr'}\,\big(r\partial_x\tilde\varphi_r\big)\big(r'\partial_x\tilde\varphi_{r'}\big),\\
\tilde\rho_r&=\frac{r}{2\pi}\,\partial_x\tilde\varphi_r,\qquad
f_{rr'}\equiv f(rk_F\,;\,r'k_F),
\end{align*}
（与 (5.3.12) 对应的哈密顿量 $H=\sum_r\int\mathrm{d}x\,\frac{2\pi^{2}v_F^{*}}{k_F}\tilde\rho_r^{2}+\frac12\sum_{rr'}\int\mathrm{d}x\,f_{rr'}\tilde\rho_r\tilde\rho_{r'}$）。无相互作用时每支是手征声子：由运动方程 $\partial_t\tilde\varphi_r=rv_F^{*}\partial_x\tilde\varphi_r$ 得 $\omega=v_F^{*}|q|$，这正是 Tomonaga--Luttinger 液体的两支线性激发。定义 $\varphi\equiv-(\tilde\varphi_{+}-\tilde\varphi_{-})$（总密度 $\tilde\rho_{+}+\tilde\rho_{-}=\frac{\partial_x\varphi}{2\pi}$ 型组合）与共轭组合 $\Theta\equiv\tilde\varphi_{+}+\tilde\varphi_{-}$，上式即标准 TLL 拉格朗日量。

\textbf{(b)} 把 $\rho=(2\pi)^{-1}\partial_x\varphi$ 代入式 (5.1.3)。用 $\partial_x(-\partial_x^{2})^{-1}\partial_x=-1$：
\begin{equation*}
L=\int\mathrm{d}x\,\Big[\frac{m}{2\rho_0}\,\frac{1}{4\pi^{2}}\dot\varphi\,(-\partial_x^{2})^{-1}\!\cdot(\!-\partial_x^{2})\,\dot\varphi
-\frac{1}{2\chi}\frac{(\partial_x\varphi)^{2}}{4\pi^{2}}\Big]
=\int\mathrm{d}x\,\Big[\pm\frac{m}{8\pi^{2}\rho_0}\dot\varphi^{2}-\frac{1}{8\pi^{2}\chi}(\partial_x\varphi)^{2}\Big],
\end{equation*}
（$\dot\rho(-\partial^{2})^{-1}\dot\rho=\mp\dot\varphi^{2}/4\pi^{2}$，符号依 $\rho=\pm\partial_x\varphi/2\pi$ 的约定；取约定使正则结构标准）。正则动量
\begin{equation*}
\tilde\varphi=\frac{\partial L}{\partial\dot\varphi}=\frac{m}{4\pi^{2}\rho_0}\dot\varphi
\qquad\Longrightarrow\qquad
\dot\varphi=\frac{4\pi^{2}\rho_0}{m}\tilde\varphi ,
\end{equation*}
\begin{equation*}
\boxed{\;H=\int\mathrm{d}x\,\Big[\frac{2\pi^{2}\rho_0}{m}\,\tilde\varphi^{2}+\frac{1}{8\pi^{2}\chi}(\partial_x\varphi)^{2}\Big],
\qquad
S=\int\mathrm{d}t\,\mathrm{d}x\,\Big[\tilde\varphi\dot\varphi-\frac{2\pi^{2}\rho_0}{m}\tilde\varphi^{2}-\frac{1}{8\pi^{2}\chi}(\partial_x\varphi)^{2}\Big]\;}
\end{equation*}
（$\rho=\partial_x\varphi/2\pi$ 的符号约定使动能项系数取正；两种约定给出相同的 $H$。）这是一对共轭标量场的二阶形式，运动方程 $\ddot\varphi=v^{2}\partial_x^{2}\varphi$，$v^{2}=\rho_0/m\chi$，正是一维流体动力学声子。

\textbf{与 (a) 的比较。}代入自由费米子的数值 $\rho_0=k_F/\pi$、$\chi=N(E_F)=1/\pi v_F$（$v=v_F$）：
\begin{equation*}
H=\int\mathrm{d}x\,\Big[2\pi v_F\,\tilde\varphi^{2}+\frac{v_F}{8\pi}(\partial_x\varphi)^{2}\Big].
\end{equation*}
而 (a) 中以 $\Theta=2k_F\partial_x^{-1}\tilde\varphi$（使 $\frac{1}{2k_F}\partial_x\varphi\,\partial_t\Theta=\tilde\varphi\dot\varphi$）改写并弃去常数后，
\begin{equation*}
H_{(a)}=\int\mathrm{d}x\,\Big[\frac{v_F}{2k_F}(\partial_x\Theta)^{2}+\frac{v_F}{2k_F}(\partial_x\varphi)^{2}\Big]
=\int\mathrm{d}x\,\Big[2\pi v_F\,\tilde\varphi^{2}+\frac{v_F}{8\pi}(\partial_x\varphi)^{2}\Big],
\end{equation*}
逐项一致（用了 $\rho_0=k_F/\pi$、$\chi=1/\pi v_F$）。所以 (b) 的作用量与 (a) 的玻色化拉格朗日量是\textbf{同一理论}：(5.1.3) 的流体动力学模型经 $\rho=\partial_x\varphi/2\pi$ 正则化后正是自由（或弱相互作用）一维费米子的 Tomonaga--Luttinger 液体；差别只在场的线性重组与归一化。这也再次印证：一维流体动力学=玻色化，且（由 (a) 的两支结构）它同时包含 $k\approx0$ 与（经紧致性，见 5.1.6）$\pm2nk_F$ 的低能激发。

% ---------------------------------------------------------------
\section*{问题 5.3.5}
\begin{problem}
均匀磁场 $\pmb{B}=B\pmb{z}$ 下，二维约化玻尔兹曼方程 (5.3.9) 中的力为 $\pmb{F}(\pmb{k},\pmb{x})=\frac{e}{c}\pmb{v}^{*}(\pmb{k})\times\pmb{B}$。(a) 证明线性化（$\tau^{-1}=0$）后方程为 $(\mathrm{i}\partial_t+\mathrm{i}\omega_c\partial_\theta-v_F^{*}q\cos(\theta-\theta_{\pmb q}))\tilde\rho-\frac{qk_F\cos(\theta_{\pmb q}-\theta)}{(2\pi)^{2}}\int f(\theta-\theta')\tilde\rho(\theta')=0$，并求 $\omega_c$；(b) 证明集体模谱由 $\tilde H_{ll'}=\omega_cl\delta_{ll'}+\frac{q}{2}(\delta_{l,l'+1}+\delta_{l,l'-1})\sqrt{v_F^{*}+\frac{k_Ff_l}{2\pi}}\sqrt{v_F^{*}+\frac{k_Ff_{l'}}{2\pi}}$ 的本征值决定，且谱关于零对称；(c) 证明 $f_l=0$ 时本征值为 $\omega_c\times$整数。
\end{problem}
\solution

\textbf{(a)}$\pmb{F}=\frac{e}{c}\pmb{v}^{*}\times\pmb{B}$：$\pmb{v}^{*}=v_F^{*}\hat{\pmb{k}}$，$\hat{\pmb{k}}\times\hat{\pmb{z}}=-\hat{\pmb{k}}_{\perp}$，故 $\pmb{F}=-\frac{ev_F^{*}B}{c}\hat{\pmb{k}}_{\perp}$。代入式 (5.3.9) 并线性化（$\partial_{\pmb{k}}\tilde\epsilon\to v_F^{*}\hat{\pmb{k}}$；$\hat{\pmb{k}}\cdot\pmb{F}=0$ 消去 $-\hat{\pmb{k}}\tilde\rho\cdot\pmb{F}$ 项；$\hat{\pmb{k}}_{\perp}\partial_\theta\tilde\rho\cdot\partial_{\pmb{x}}\tilde\epsilon$ 为二阶舍去），力项只剩
\begin{equation*}
k_F^{-1}\,\hat{\pmb{k}}_{\perp}\!\cdot\!\pmb{F}\;\partial_\theta\tilde\rho
=-\frac{ev_F^{*}B}{c\,k_F}\,\partial_\theta\tilde\rho
=-\omega_c\,\partial_\theta\tilde\rho,
\qquad
\boxed{\;\omega_c=\frac{eB}{m^{*}c}\;}
\end{equation*}
（用 $m^{*}v_F^{*}=k_F$）。于是
\begin{equation*}
\partial_t\tilde\rho+v_F^{*}\hat{\pmb{k}}\cdot\partial_{\pmb{x}}\tilde\rho
+\frac{k_F}{(2\pi)^{2}}\int\mathrm{d}\theta'\,f(\theta-\theta')\,\hat{\pmb{k}}\cdot\partial_{\pmb{x}}\tilde\rho(\theta')
-\omega_c\partial_\theta\tilde\rho=0 .
\end{equation*}
对 $\pmb{x}$ 傅里叶（$\partial_{\pmb{x}}\to\mathrm{i}\pmb{q}$，$\hat{\pmb{k}}\cdot\pmb{q}=q\cos(\theta_{\pmb{q}}-\theta)$）并乘 $\mathrm{i}$：
\begin{equation*}
\Big(\mathrm{i}\partial_t+\mathrm{i}\omega_c\partial_\theta-v_F^{*}q\cos(\theta-\theta_{\pmb{q}})\Big)\tilde\rho
-\frac{qk_F\cos(\theta_{\pmb{q}}-\theta)}{(2\pi)^{2}}\int\mathrm{d}\theta'\,f(\theta-\theta')\,\tilde\rho(\theta',\pmb{q},t)=0 ,
\end{equation*}
即所求（注意 $\cos$ 为偶函数，两个 $\cos$ 写法等价）。

\textbf{(b)}取 $\theta_{\pmb q}=0$，展开 $\tilde\rho=\sum_l\tilde\rho_le^{\mathrm{i}l\theta}/\sqrt{2\pi}$、$f=\sum_lf_le^{\mathrm{i}l\theta}$，用 $\cos\theta\to\frac12(\delta_{l,l\pm1})$ 移位、$\partial_\theta\to\mathrm{i}l$、$f_l=f_{-l}\in\mathbb{R}$：
\begin{equation*}
\mathrm{i}\partial_t\tilde\rho_l
=\omega_c l\,\tilde\rho_l
+\frac{q}{2}\Big(\sqrt{v_F^{*}+\frac{k_Ff_{l-1}}{2\pi}}\sqrt{v_F^{*}+\frac{k_Ff_{l}}{2\pi}}\,\tilde\rho_{l-1}
+\sqrt{v_F^{*}+\frac{k_Ff_{l}}{2\pi}}\sqrt{v_F^{*}+\frac{k_Ff_{l+1}}{2\pi}}\,\tilde\rho_{l+1}\Big),
\end{equation*}
即 $\mathrm{i}\partial_t\tilde\rho=\tilde H\tilde\rho$，其中
\begin{equation*}
\tilde H_{ll'}=\omega_cl\,\delta_{ll'}
+\frac{q}{2}\big(\delta_{l,l'+1}+\delta_{l,l'-1}\big)
\sqrt{v_F^{*}+\frac{k_Ff_l}{2\pi}}\sqrt{v_F^{*}+\frac{k_Ff_{l'}}{2\pi}} ,
\end{equation*}
$\tilde H$ 实对称，集体模能量为其本征值。\textbf{关于零对称}：定义反线性映射 $(A\psi)_l=(-1)^{l}\psi^{*}_{-l}$。用 $f_l=f_{-l}$（故 $\sqrt{\cdot}$ 为偶）与实性，
\begin{align*}
(\tilde HA\psi)_l
&=\omega_cl(-1)^{l}\psi^{*}_{-l}
+\frac{q}{2}\big[(-1)^{l-1}\sqrt{m_{l-1}}\sqrt{m_l}\,\psi^{*}_{-(l-1)}
+(-1)^{l+1}\sqrt{m_{l+1}}\sqrt{m_l}\,\psi^{*}_{-(l+1)}\big]\\
&=(-1)^{l}\Big[\omega_cl\psi^{*}_{-l}
-\frac{q}{2}\big(\sqrt{m_{l-1}}\sqrt{m_l}\,\psi^{*}_{1-l}+\sqrt{m_{l+1}}\sqrt{m_l}\,\psi^{*}_{-l-1}\big)\Big],
\end{align*}
（$m_l\equiv v_F^{*}+\frac{k_Ff_l}{2\pi}$，$m_{l\mp1}=m_{\mp(l\mp1)}$ 偶）。若 $\tilde H\psi=E\psi$，在指标 $-l$ 处取方程并取复共轭：$E\psi^{*}_{-l}=-\omega_cl\psi^{*}_{-l}+\frac{q}{2}(\sqrt{m_{l+1}}\psi^{*}_{-l-1}+\sqrt{m_{l-1}}\psi^{*}_{1-l})\sqrt{m_l}$，于是上式中括号 $=-E\psi^{*}_{-l}$，即 $\tilde H(A\psi)=-E(A\psi)$：本征值成对 $\pm E$（$E=0$ 自映射），\textbf{谱关于零对称}。物理上正频率对应玻色模式的产生、负频率对应湮灭（见 5.3.3 节脚注）。

\textbf{(c)}$f_l=0$ 时在 $\theta$ 表象中 $\tilde H=-\mathrm{i}\omega_c\partial_\theta+v_F^{*}q\cos(\theta-\theta_{\pmb q})$（$\partial_\theta\to$ 与 $-\mathrm{i}\partial_\theta$ 对应）。作规范（幺正）变换 $U=\exp\!\big[-\mathrm{i}\frac{v_F^{*}q}{\omega_c}\sin(\theta-\theta_{\pmb q})\big]$：因 $U^{-1}(-\mathrm{i}\partial_\theta)U=-\mathrm{i}\partial_\theta-\frac{v_F^{*}q}{\omega_c}\cos(\theta-\theta_{\pmb q})$，
\begin{equation*}
\tilde H=U\big(-\mathrm{i}\omega_c\partial_\theta\big)U^{-1} .
\end{equation*}
而 $-\mathrm{i}\omega_c\partial_\theta$ 的谱为 $\{\omega_cl\,|\,l\in\mathbb{Z}\}$（本征函数 $\mathrm{e}^{\mathrm{i}l\theta}$），故
\begin{equation*}
\boxed{\;\mathrm{spec}(\tilde H)=\{\,\omega_c\times\text{整数}\,\},\qquad
\tilde\rho\propto\exp\big(\mathrm{i}l\theta-\mathrm{i}\tfrac{v_F^{*}q}{\omega_c}\sin(\theta-\theta_{\pmb q})\big)\;}
\end{equation*}
（即提示的试探函数，本征值 $\omega_cl$）。这样流体动力学理论精确恢复了自由费米子在磁场中的结果：粒子--空穴激发能量为 Landau 能级间距 $\omega_c$ 的整数倍。

% ---------------------------------------------------------------
\section*{问题 5.3.6}
\begin{problem}
二维费米液体的交流输运：$\pmb{F}=e\pmb{E}\,\mathrm{e}^{\mathrm{i}\omega t-\pmb{q}\cdot\pmb{x}}$。(a) $f=0$ 时求 $\tilde\rho$；(b) 准确到 $f$ 一阶求 $\tilde\rho$；(c) 求感应电流 $e\pmb{J}=\sigma(\omega,\pmb{q})\pmb{E}$ 与光学电导 $\sigma(\omega,\pmb{q})$（准确到 $f$ 一阶）。
\end{problem}
\solution

记 $D(\theta)\equiv\tau^{-1}+\mathrm{i}\omega-\mathrm{i}v_F^{*}\hat{\pmb{k}}\cdot\pmb{q}$，$b\equiv v_F^{*}q$，$a\equiv\tau^{-1}+\mathrm{i}\omega$。取 $\pmb{E}\parallel\pmb{q}$ 沿 $x$（横向分量独立且由各向同性解耦），$\hat{\pmb{k}}\cdot\pmb{E}=E\cos\theta$。

\textbf{(a)}$f=0$：式 (5.3.10) 傅里叶后为 $(\tau^{-1}+\mathrm{i}\omega-\mathrm{i}v_F^{*}\hat{\pmb{k}}\cdot\pmb{q})\tilde\rho=\frac{k_Fe}{(2\pi)^{2}}\hat{\pmb{k}}\cdot\pmb{E}$，故
\begin{equation*}
\boxed{\;\tilde\rho^{(0)}(\theta)=\frac{k_FeE}{(2\pi)^{2}}\cdot\frac{\cos\theta}{\tau^{-1}+\mathrm{i}\omega-\mathrm{i}v_F^{*}q\cos\theta}\;}
\end{equation*}

\textbf{(b)}保留 $\int f\partial_{\pmb{x}}\tilde\rho$ 项到一阶：$(\tau^{-1}+\mathrm{i}\omega-\mathrm{i}v_F^{*}\hat{\pmb{k}}\cdot\pmb{q})\tilde\rho-\frac{k_Fe}{(2\pi)^{2}}\hat{\pmb{k}}\cdot\pmb{E}+\frac{k_F}{(2\pi)^{2}}\hat{\pmb{k}}\cdot(\mathrm{i}\pmb{q})\int f(\theta-\theta')\tilde\rho(\theta')=0$ 迭代一次：
\begin{equation*}
\boxed{\;\tilde\rho(\theta)=\tilde\rho^{(0)}(\theta)
+\mathrm{i}\,\frac{k_Fq\cos\theta}{(2\pi)^{2}D(\theta)}
\int\mathrm{d}\theta'\,f(\theta-\theta')\,\tilde\rho^{(0)}(\theta')\;}
\end{equation*}
（对 $f=f_1$ 分量：$\int f(\theta-\theta')\cos\theta'\,\mathrm{d}\theta'=2\pi f_1\cos\theta$，代入即显式解。）

\textbf{(c)}由式 (5.3.6)，$\tilde{\pmb{v}}=\hat{\pmb{k}}\big[v_F^{*}+\frac{k_Ff_1}{\pi}\big]$（只有 $l=\pm1$ 分量进入回流：$\int\frac{\mathrm{d}^{2}k'}{(2\pi)^{2}}fv^{*}\delta(\xi^{*})=\hat{\pmb{k}}\frac{k_Ff_1}{\pi}$），
\begin{equation*}
e\pmb{J}=e\int\mathrm{d}\theta\,\tilde\rho(\theta)\,\hat{\pmb{k}}\Big(v_F^{*}+\frac{k_Ff_1}{\pi}\Big)
\equiv\sigma(\omega,\pmb{q})\,\pmb{E},
\end{equation*}
准确到 $f$ 一阶：
\begin{equation*}
\sigma=e\Big(v_F^{*}\!\int\!\tilde\rho^{(0)}\cos\theta\,\mathrm{d}\theta\Big)
+e\Big(v_F^{*}\!\int\!\tilde\rho^{(1)}\cos\theta\,\mathrm{d}\theta\Big)
+\frac{e\,k_Ff_1}{\pi}\!\int\!\tilde\rho^{(0)}\cos\theta\,\mathrm{d}\theta .
\end{equation*}
第一项（精确的 $\omega,q$ 依赖）：用 $\int_0^{2\pi}\frac{\cos^{2}\theta\,\mathrm{d}\theta}{a-\mathrm{i}b\cos\theta}=\frac{2\pi}{b^{2}}\Big[\frac{a^{2}}{\sqrt{a^{2}-b^{2}}}-a\Big]$，
\begin{equation*}
\boxed{\;\sigma^{(0)}(\omega,q)=\frac{e^{2}k_Fv_F^{*}}{(2\pi)^{2}}\int\frac{\cos^{2}\theta\,\mathrm{d}\theta}{a-\mathrm{i}b\cos\theta}
=\frac{e^{2}k_F}{2\pi qv_F^{*}}\Big[\frac{a^{2}}{\sqrt{a^{2}-b^{2}}}-a\Big]\;}
\end{equation*}
第三项立即给出回流修正 $\sigma^{(3)}=\sigma^{(0)}\cdot\frac{k_Ff_1}{\pi v_F^{*}}$；第二项
\begin{equation*}
\sigma^{(2)}=\mathrm{i}\,\frac{e^{2}k_F^{2}v_F^{*}q}{(2\pi)^{4}}\!\int\!\mathrm{d}\theta\,\mathrm{d}\theta'\,
\frac{f(\theta-\theta')\,\cos^{2}\theta\cos\theta'}{D(\theta)D(\theta')}
=\mathrm{i}\,\frac{e^{2}k_F^{2}v_F^{*}q\,f_1}{(2\pi)^{3}}\;\Big|\!\int\!\frac{\cos\theta\,\mathrm{e}^{\mathrm{i}\theta}\,\mathrm{d}\theta}{a-\mathrm{i}b\cos\theta}\Big|^{2},
\end{equation*}
（末等式用 $f_1=f_{-1}$ 与留数公式 $\int\frac{\mathrm{e}^{\mathrm{i}n\theta}\mathrm{d}\theta}{a-b\cos\theta}=\frac{2\pi}{\sqrt{a^{2}-b^{2}}}\big(\frac{a-\sqrt{a^{2}-b^{2}}}{b}\big)^{n}$）。故
\begin{equation*}
\boxed{\;\sigma(\omega,q)=\sigma^{(0)}(\omega,q)\Big[1+\frac{k_Ff_1}{\pi v_F^{*}}\Big]+\sigma^{(2)}(\omega,q)\;}
\end{equation*}
检验极限：
(i) $\sigma(\omega,0)$：$q\to0$ 时 $\tilde\rho^{(1)}\propto q$ 消失，得 Drude 型光学电导
\begin{equation*}
\sigma(\omega,0)=\frac{e^{2}\rho}{k_F}\cdot\frac{v_F^{*}+\frac{k_Ff_1}{\pi}}{\tau^{-1}+\mathrm{i}\omega}
=\sigma_{\mathrm{Drude}}(\omega)\cdot\Big(1+\frac{k_Ff_1}{\pi v_F^{*}}\Big),
\qquad
\sigma_{\mathrm{Drude}}(\omega)=\frac{e^{2}\tau\rho/m^{*}}{1+\mathrm{i}\omega\tau},
\end{equation*}
即 Drude 形式乘以回流增强因子 $|\tilde{\pmb v}|/v_F^{*}=1+\frac{k_Ff_1}{\pi v_F^{*}}$，与正文 $\sigma=\frac{e^{2}\tau|\tilde{\pmb v}|\rho}{k_F}$ 一致。
(ii) $\omega=0$、$\tau\to\infty$（干净极限、有限 $q$）：$a\to0$，$\sigma^{(0)}\to0$——无碰撞时有限波矢的直流输运无耗散，合理。

% ---------------------------------------------------------------
\section*{问题 5.4.1}
\begin{problem}
(a) 计算 $Z/Z_0$ 微扰展开中 $V^{2}$ 量级的项并画出费曼图；(b) 计算 $\ln(Z/Z_0)$ 中 $V^{2}$ 量级的项并画图。
\end{problem}
\solution

记一阶插人算符 $O\equiv\int\mathrm{d}x\,\mathrm{d}x'\,\psi^{*}(x)\psi(x)[-\mathrm{i}V(x,x')]\psi^{*}(x')\psi(x')$，则 $Z/Z_0=1+\langle O\rangle_0+\frac{1}{2}\langle O^{2}\rangle_0+\cdots$，其中 $O^{2}$ 含两个相互作用积分 $\int\mathrm{d}1\,\mathrm{d}2\,(\cdots V(1,2)\text{ 两次})$。

\textbf{(a)} 用 Wick 定理展开 $\frac12\langle O^{2}\rangle_0$，按两顶点之间是否有费米线相连分成两类：

\emph{非连通部分}：两顶点各自内部收缩，不给 $\frac12\langle O\rangle_0^{2}$。它对应三张由两个一阶图（图 5.7 的 Hartree 与 Fock 图）拼成的非连通图：Hartree$+$Hartree、Hartree$+$Fock、Fock$+$Fock。

\emph{连通部分} $\frac12\langle O^{2}\rangle_{0c}$：两顶点由费米线相连，只有两种拓扑：

（i）\textbf{环链（二阶直接）图}：两条虚线都连接顶点 1、2，费米线构成两个独立的小环（每个环穿过两顶点各一次）：
\begin{equation*}
D_{2}=\frac{1}{2}\int\mathrm{d}1\,\mathrm{d}2\;(-1)^{2}\,
\big[\mathrm{i}G_0(1,2)\big]^{2}\big[\mathrm{i}G_0(2,1)\big]^{2}\,
[-\mathrm{i}V(1,2)]\,[-\mathrm{i}V(2,1)],
\end{equation*}
对称因子 $1/2$；这是 RPA 气泡链级数 $-\sum_n\frac1n(\mathrm{i}\Pi^{00}[-\mathrm{i}V])^{n}$ 的 $n=2$ 项。

（ii）\textbf{交换（二阶 Fock）图}：费米线拧成一条穿过两顶点各两次的单圈（两虚线交叉），多一个闭合圈因子 $(-1)$：
\begin{equation*}
E_{2}=-\frac{1}{2}\int\mathrm{d}1\,\mathrm{d}2\;2\,\big[\mathrm{i}G_0(1,2)\big]^{4}\,
[-\mathrm{i}V(1,2)]\,[-\mathrm{i}V(1,2)]^{*}\cdot(\text{排列因子})
=-\int\mathrm{d}1\,\mathrm{d}2\,\big[\mathrm{i}G_0(1,2)\big]^{4}[-\mathrm{i}V(1,2)]^{2},
\end{equation*}
（对称因子 1；圈数 1 给 $-1$。）

于是
\begin{equation*}
\boxed{\;
\frac{Z}{Z_0}\Big|_{V^{2}}=\frac{1}{2}\langle O\rangle_0^{2}+D_{2}+E_{2}\;}
\end{equation*}
费曼图：$\frac12\langle O\rangle_0^{2}$ 为三张非连通双图；$D_2$ 为双气泡链图（两圈两虚线）；$E_2$ 为交叉双虚线单圈图。

\textbf{(b)} 对 $\ln(Z/Z_0)$ 用 $\ln(1+x)=x-\frac{x^{2}}{2}+\cdots$，$x=\langle O\rangle_0+\frac12\langle O^{2}\rangle_0$：
\begin{equation*}
\ln\frac{Z}{Z_0}=\langle O\rangle_0+\frac12\langle O^{2}\rangle_0-\frac12\langle O\rangle_0^{2}+O(V^{3})
=\langle O\rangle_0+\frac12\langle O^{2}\rangle_{0c}+O(V^{3}),
\end{equation*}
即非连通的 $\frac12\langle O\rangle_0^{2}$ 恰好被消去（这正是连链定理 (5.4.2)）：
\begin{equation*}
\boxed{\;\ln\frac{Z}{Z_0}\Big|_{V^{2}}=D_{2}+E_{2}\;}
\end{equation*}
只剩两张连通图（环链图与二阶交换图）。物理解释：$Z/Z_0$ 中的非连通图是独立热涨落的乘积，取对数后（自由能是广延量）只保留连通图。

% ---------------------------------------------------------------
\section*{问题 5.4.2}
\begin{problem}
用微扰论证明式 (5.4.4)。注意 $\int\frac{\mathrm{d}\nu}{2\pi}\mathrm{i}G_{0,\pmb{k}'\nu}\mathrm{e}^{\mathrm{i}0^{+}\nu}$ 代表等时 $\langle c^{\dagger}c\rangle$，$\int\frac{\mathrm{d}\nu}{2\pi}\mathrm{i}G_{0,\pmb{k}'\nu}\mathrm{e}^{-\mathrm{i}0^{+}\nu}$ 代表等时 $\langle cc^{\dagger}\rangle$。
\end{problem}
\solution

由式 (5.4.3)，准确到 $V$ 一阶：
\begin{equation*}
\mathrm{i}G_{\pmb{k}\omega}=\mathrm{i}G_{0,\pmb{k}\omega}
+\big\langle c_{\pmb{k}}(t)c_{\pmb{k}}^{\dagger}(0)\,(-\mathrm{i})\!\int\!\mathrm{d}t'\,H_I(t')\big\rangle_{0c} .
\end{equation*}
一阶项用 Wick 定理收缩，拓扑上有两个图（图 5.8）：

\textbf{Hartree 图（图 5.8(a)）}：外线 $c\to c^{\dagger}$ 直通，顶点处剩下一对同点算符缩成闭合小圈。数值为
\begin{equation*}
(-1)\,(\mathrm{i}G_{0,\pmb{k}\omega})^{2}(-\mathrm{i}V_0)\,
\frac{1}{\mathcal V}\sum_{\pmb{k}'\alpha'}\int\frac{\mathrm{d}\nu}{2\pi}\,\mathrm{i}G_{0,\pmb{k}'\nu}\,\mathrm{e}^{\mathrm{i}0^{+}\nu},
\end{equation*}
其中 $(-1)$ 来自闭合费米圈。\textbf{频率积分}：
\begin{equation*}
\int\frac{\mathrm{d}\nu}{2\pi}\,\frac{\mathrm{i}\,\mathrm{e}^{\mathrm{i}0^{+}\nu}}{\nu-\xi_{\pmb{k}'}+\mathrm{i}0^{+}\operatorname{sgn}\nu}:
\end{equation*}
被积函数的极点在 $\nu=\xi_{\pmb{k}'}-\mathrm{i}0^{+}$（来自 $\nu>0$ 支）与 $\nu=\xi_{\pmb{k}'}+\mathrm{i}0^{+}$（来自 $\nu<0$ 支）；因子 $\mathrm{e}^{\mathrm{i}0^{+}\nu}=\mathrm{e}^{\mathrm{i}0^{+}\operatorname{Re}\nu}\mathrm{e}^{-0^{+}\operatorname{Im}\nu}$ 要求在上半平面闭合，故只有 $\nu<0$ 支的极点（位于上半平面、且仅当 $\xi_{\pmb{k}'}<0$ 时在实轴附近被俘获）有贡献：
\begin{equation*}
\int\frac{\mathrm{d}\nu}{2\pi}\,\mathrm{i}G_{0,\pmb{k}'\nu}\mathrm{e}^{\mathrm{i}0^{+}\nu}
=-\Theta(-\xi_{\pmb{k}'})=-n_{\pmb{k}'}\equiv-\langle c^{\dagger}_{\pmb{k}'}c_{\pmb{k}'}\rangle ,
\end{equation*}
即等时的 $\langle c^{\dagger}c\rangle$（带一圈负号）。于是 Hartree 项 $=(-1)(\mathrm{i}G_0)^{2}(-\mathrm{i}V_0)\cdot\frac{1}{\mathcal V}\sum(-n_{\pmb{k}'})$。

\textbf{Fock 图（图 5.8(b)）}：外线经过顶点（交换式收缩），顶点处剩下等时对 $\langle cc^{\dagger}\rangle$：
\begin{equation*}
(\mathrm{i}G_{0,\pmb{k}\omega})^{2}\,\frac{1}{\mathcal V}\sum_{\pmb{q}}\int\frac{\mathrm{d}\nu}{2\pi}\,
\mathrm{i}G_{0,\pmb{k}-\pmb{q},\nu}\,(-\mathrm{i}V_{\pmb{q}})\,\mathrm{e}^{-\mathrm{i}0^{+}\nu},
\end{equation*}
此时 $\mathrm{e}^{-\mathrm{i}0^{+}\nu}$ 要求下半平面闭合，俘获 $\nu>0$ 支的极点（仅当 $\xi_{\pmb{k}-\pmb{q}}>0$）：
\begin{equation*}
\int\frac{\mathrm{d}\nu}{2\pi}\,\mathrm{i}G_{0,\pmb{k}-\pmb{q},\nu}\mathrm{e}^{-\mathrm{i}0^{+}\nu}
=+\Theta(\xi_{\pmb{k}-\pmb{q}})=1-n_{\pmb{k}-\pmb{q}}=\langle c_{\pmb{k}-\pmb{q}}c^{\dagger}_{\pmb{k}-\pmb{q}}\rangle .
\end{equation*}

合并两图即得式 (5.4.4)：
\begin{equation*}
\boxed{\;
\mathrm{i}G_{\pmb{k}\omega}=\mathrm{i}G_{0,\pmb{k}\omega}
+(-1)(\mathrm{i}G_{0,\pmb{k}\omega})^{2}(-\mathrm{i}V_0)\frac{1}{\mathcal V}\sum_{\pmb{k}'\alpha'}\int\frac{\mathrm{d}\nu}{2\pi}\,\mathrm{i}G_{0,\pmb{k}'\nu}\mathrm{e}^{\mathrm{i}0^{+}\nu}
+(\mathrm{i}G_{0,\pmb{k}\omega})^{2}\frac{1}{\mathcal V}\sum_{\pmb{q}}\int\frac{\mathrm{d}\nu}{2\pi}\,\mathrm{i}G_{0,\pmb{k}-\pmb{q},\nu}(-\mathrm{i}V_{\pmb{q}})\mathrm{e}^{-\mathrm{i}0^{+}\nu}\;}
\end{equation*}
写成 $\mathrm{i}G=\mathrm{i}G_0+(\mathrm{i}G_0)^{2}(-\mathrm{i}\Sigma)$ 后，
\begin{equation*}
\Sigma_{\pmb{k}\omega}=V_0\rho_0-\frac{1}{\mathcal V}\sum_{\pmb{q}}V_{\pmb{q}}\,n_{\pmb{k}-\pmb{q}}
\;(+\text{可并入化学势的常数}\ \mathcal V^{-1}\textstyle\sum_{\pmb{q}}V_{\pmb{q}}),
\end{equation*}
即 Hartree 贡献 $V_0\rho_0$ 加 Fock 贡献 $-\mathcal V^{-1}\sum_{\pmb{q}}V_{\pmb{q}}n_{\pmb{k}-\pmb{q}}$，与 Hartree--Fock 结果 (5.2.4)、(5.3.2) 一致。

% ---------------------------------------------------------------
\section*{问题 5.4.3}
\begin{problem}
写出图 5.11(a)--(e) 所示 $V_{\pmb q}^{2}$ 量级自能各项的表达式（图 5.11(f) 已含于 Dyson 求和，不计）。
\end{problem}
\solution

记 $\displaystyle\int_{\pmb q}\equiv\frac{1}{\mathcal V}\sum_{\pmb q}\int\frac{\mathrm{d}\nu}{2\pi}$，$\mathrm{i}G_{\pmb q\nu}=\mathrm{i}/(\nu-\xi_{\pmb q})$，每个闭合费米圈因子 $(-1)$。五张图对 $-\mathrm{i}\Sigma_{\pmb{k}\omega}$ 的贡献如下（每项再乘外线因子 $(\mathrm{i}G_{0,\pmb{k}\omega})^{2}$ 进入 $\mathrm{i}G$）：

\textbf{(a) 交换修饰的蝌蚪图}（Hartree 图中密度圈被一条交换虚线修饰）：外虚线携带 $V_0$（零动量转移），圈上的虚线携带 $(\pmb q,\nu)$：
\begin{equation*}
-\mathrm{i}\Sigma^{(a)}
=(-\mathrm{i}V_0)(-1)\int_{\pmb q}(-\mathrm{i}V_{\pmb q})(-1)
\int_{\pmb p}\,\mathrm{i}G_{0,\pmb p\xi}\,\mathrm{i}G_{0,\pmb p-\pmb q,\xi-\nu}\,\mathrm{e}^{\mathrm{i}0^{+}\xi} ,
\end{equation*}
即 $V_0$ 乘以密度圈的一阶（交换型）修正。物理：Hartree 项中背景密度的 $V^{2}$ 修整。

\textbf{(b) 直接修饰的蝌蚪图}（两个费米圈由一条虚线相连）：外虚线仍携带 $V_0$，圈间虚线携带 $(\pmb q,\nu)$，中间是一个气泡 $\mathrm{i}\Pi^{00}$：
\begin{equation*}
-\mathrm{i}\Sigma^{(b)}
=(-\mathrm{i}V_0)(-1)\int_{\pmb q}(-\mathrm{i}V_{\pmb q})\;\mathrm{i}\Pi^{00}_{\pmb q\nu}\,,
\qquad
\mathrm{i}\Pi^{00}_{\pmb q\nu}=(-1)\int_{\pmb p}\,\mathrm{i}G_{0,\pmb p\xi}\,\mathrm{i}G_{0,\pmb p+\pmb q,\xi+\nu} ;
\end{equation*}
这是 RPA 屏蔽级数播种在 Hartree 道上的 $V^{2}$ 项（$V_0\,\Pi^{00}V$ 型）。

\textbf{(c) 二阶交换图}（外线上两条交叉虚线；无闭合圈）：
\begin{equation*}
-\mathrm{i}\Sigma^{(c)}
=(-1)\int_{\pmb q}\!\int_{\pmb p}(-\mathrm{i}V_{\pmb q})(-\mathrm{i}V_{\pmb p})\,
\mathrm{i}G_{0,\pmb k-\pmb q,\omega-\nu}\,
\mathrm{i}G_{0,\pmb k-\pmb q-\pmb p,\omega-\nu-\xi}\,
\mathrm{i}G_{0,\pmb k-\pmb p,\omega-\xi} ,
\end{equation*}
额外的 $(-1)$ 来自交叉收缩（费米线的一次置换）。

\textbf{(d) 气泡挂在双虚线上的图}（外线两点经两条虚线与一个费米圈相连；Fock 道的屏蔽种子）：
\begin{equation*}
-\mathrm{i}\Sigma^{(d)}
=(-1)\int_{\pmb q}(-\mathrm{i}V_{\pmb q})\,\mathrm{i}\Pi^{00}_{\pmb q\nu}\,(-\mathrm{i}V_{\pmb q})\,
\mathrm{i}G_{0,\pmb k-\pmb q,\omega-\nu} ,
\end{equation*}
即 $G\times[V\Pi^{00}V]$：Fock 图中相互作用被一阶粒子--空穴激发 dressing。

\textbf{(e) 二阶直接图}（外线上两条嵌套（不交叉）虚线；无闭合圈）：
\begin{equation*}
-\mathrm{i}\Sigma^{(e)}
=\int_{\pmb q}\!\int_{\pmb p}(-\mathrm{i}V_{\pmb q})(-\mathrm{i}V_{\pmb p})\,
\mathrm{i}G_{0,\pmb k-\pmb q,\omega-\nu}\,
\mathrm{i}G_{0,\pmb k-\pmb q-\pmb p,\omega-\nu-\xi}\,
\mathrm{i}G_{0,\pmb k-\pmb p,\omega-\xi} ,
\end{equation*}
与 (c) 具有相同的 $G$--$V$ 骨架，差别只在符号（无交叉因子）与收敛因子 $\mathrm{e}^{\pm\mathrm{i}0^{+}\nu}$ 的分配（对应不同的密度简并统计因子）。

\textbf{(f)} 为 $\Sigma^{(1)}G_0\Sigma^{(1)}$（两个一阶自能串在一起）：它已在图 5.10 的 Dyson 求和 $\mathrm{i}G_0\sum_n[(\mathrm{i}G_0)(-\mathrm{i}\Sigma)]^{n}$ 中被计入，若再算入即重复计数，故对自能无贡献。

以上 (a)--(e) 正是 $V^{2}$ 量级全部的一粒子不可约（1PI）自能图：把 $\Sigma_{\pmb{k}\omega}=\mathrm{i}(\text{各图之和})/\mathrm{i}$ 代入 $G_{\pmb{k}\omega}=1/(\omega-\xi_{\pmb{k}}-\Sigma_{\pmb{k}\omega}+\mathrm{i}\operatorname{sgn}(\omega)0^{+})$ 即得含屏蔽与衰变效应的准粒子性质（例如 (d)+(e) 的虚部给出发射率 $\mathrm{Im}\Sigma\propto\xi^{2}$ 的准粒子衰变，见 5.4.4 节）。

% ---------------------------------------------------------------
\section*{问题 5.6.1}
\begin{problem}
证明式 (5.6.6) 与 (5.6.7)：非线性 $\sigma$ 模型系数 $g^{-1}=-a^{-d}N_{\mathrm{site}}^{-1}\partial^{2}_{B_y}E_0(B_y)|_0$、$v^{2}/g=a^{-d}\partial^{2}_{\pmb a}E_0(\pmb a)|_0$，其中 $E_0(B_y)=N_{\mathrm{site}}^{-1}\sum'_{\pmb k}(-E^{g}_{+,\pmb k}-E^{g}_{-,\pmb k})$，$E^{g}_{\pm,\pmb k}=\sqrt{(\xi_{\pmb k}\pm B_y/2)^{2}+B_0^{2}}$；$E^{v}_{\pm,\pmb k}=\big|\sqrt{\frac14(\xi_{\pmb k+\pmb a/2}+\xi_{\pmb k-\pmb a/2})^{2}+B_0^{2}}\pm\frac12(\xi_{\pmb k+\pmb a/2}-\xi_{\pmb k-\pmb a/2})\big|$。
\end{problem}
\solution

\textbf{统一的对角化。}在约化布里渊区基 $\psi_{\pmb k}=(c_{\pmb k},c_{\pmb k+\pmb Q})$（$\otimes$ 自旋）中，均匀涨落场对应的平均场哈密顿量都取 $4\times4$ 形式
\begin{equation*}
H=\begin{pmatrix}
\xi_{\pmb k}\,\mathbb 1+\hat A & -B_0\sigma^{z}\\[2pt]
-B_0\sigma^{z} & -\xi_{\pmb k}\,\mathbb 1+\hat A
\end{pmatrix},
\end{equation*}
其中 (i) 均匀 $a_0$ 场：$\hat A=-\frac{B_y}{2}\sigma^{y}$（即 $a_0(\pmb i)=-\frac12B_y\sigma^{y}$，$\sigma^{y}$ 实使 $\phi$ 取实）；(ii) 均匀 $\pmb a$ 场：跳跃相位 $\mathrm{e}^{\mathrm{i}\pmb a\cdot\pmb\delta\,\sigma^{y}/2}$ 使 $\cos(\pmb k\cdot\pmb\delta)\to\cos(\pmb k\cdot\pmb\delta+\pmb a\cdot\pmb\delta\,\sigma^{y}/2)$，即 $\xi_{\pmb k}\to\bar\xi_{\pmb k}+\delta\xi_{\pmb k}\sigma^{y}$，$\bar\xi_{\pmb k}=\frac12(\xi_{\pmb k+\pmb a/2}+\xi_{\pmb k-\pmb a/2})$，$\delta\xi_{\pmb k}=\frac12(\xi_{\pmb k+\pmb a/2}-\xi_{\pmb k-\pmb a/2})$，$\hat A=\delta\xi_{\pmb k}\sigma^{y}$。

计算 $H^{2}$。对角块：$(\xi+\hat A)^{2}+B_0^{2}(\sigma^{z})^{2}=\xi^{2}+B_0^{2}+\hat A^{2}+\xi\{\hat A,\cdot\}\cdot\,,$ 精确地
\begin{equation*}
H^{2}=\begin{pmatrix}
\xi^{2}+B_0^{2}+\hat A^{2}+\xi\hat A+\hat A\xi & 0\\
0 & \xi^{2}+B_0^{2}+\hat A^{2}-\xi\hat A-\hat A\xi
\end{pmatrix},
\end{equation*}
（非对角块 $\propto\{(\xi+\hat A),\sigma^{z}\}$ 型组合中 $\{\sigma^{y},\sigma^{z}\}=0$ 使含 $B_0$ 的交叉项相消：$(-B_0\sigma^{z})$ 与对角块的反对易子给 $\propto\{\hat A,\sigma^{z}\}+\{\xi,\sigma^{z}\}=B_0^{-1}\cdot0$——逐项：$H^2_{12}=(\xi+\hat A)(-B_0\sigma^{z})+(-B_0\sigma^{z})(-\xi+\hat A)=-B_0\xi\sigma^{z}-B_0\hat A\sigma^{z}+B_0\xi\sigma^{z}-B_0\sigma^{z}\hat A=-B_0(\hat A\sigma^{z}+\sigma^{z}\hat A)=-B_0\{\hat A,\sigma^{z}\}=0$，因 $\hat A\propto\sigma^{y}$。）而 $\hat A^{2}=\hat a^{2}\,\mathbb 1$，$\xi\hat A+\hat A\xi=2\xi\hat a\,\sigma^{y}$（$\hat a$ 为 $\hat A$ 的系数）。故
\begin{equation*}
E^{2}=\xi^{2}+B_0^{2}+\hat a^{2}\pm\xi\hat a\cdot(\pm1)
=(\xi\pm\hat a)^{2}+B_0^{2},
\end{equation*}
四个本征值为 $\pm\sqrt{(\xi\pm\hat a)^{2}+B_0^{2}}$。

\textbf{(5.6.6)：}取 $\hat a=-B_y/2$（符号不影响 $E^{2}$）：
\begin{equation*}
E^{g}_{\pm,\pmb k}=\sqrt{\Big(\xi_{\pmb k}\pm\frac{B_y}{2}\Big)^{2}+B_0^{2}} ,
\qquad
E_0(B_y)=N_{\mathrm{site}}^{-1}\sum'_{\pmb k}\big(-E^{g}_{+,\pmb k}-E^{g}_{-,\pmb k}\big),
\end{equation*}
即式 (5.6.6)。又 $L_{\mathrm{eff}}$ 中 $a_0^{2}$ 项来自作用量 $-\!\int\!\mathrm{d}t\,E_0(B_y(t))$ 的展开：$-E_0(B_y)=-E_0(0)-\frac{B_y^{2}}{2}\partial^{2}_{B_y}E_0\big|_0$，而 $B_y=\partial_t\phi$、$\int\mathrm{d}^{d}x\,\frac{(\partial_t\pmb n)^{2}}{2g}=\int\mathrm{d}t\,\mathrm{d}^{d}x\,\frac{\dot\phi^{2}}{2g}$（$\pmb n$ 涨落即 $\phi$），比较系数（每格点能量对每格点作用量）：
\begin{equation*}
\frac{1}{2g}=-\frac{a^{-d}}{2N_{\mathrm{site}}}\partial^{2}_{B_y}E_0\big|_{0}
\quad\Longrightarrow\quad
\boxed{\;g^{-1}=-a^{-d}N_{\mathrm{site}}^{-1}\frac{\partial^{2}E_0(B_y)}{\partial B_y^{2}}\Big|_{B_y=0}\;}
\end{equation*}
（显式：$\partial^{2}_{B_y}E^{g}_{\pm}|_0=\frac{B_0^{2}}{4(\xi_{\pmb k}^{2}+B_0^{2})^{3/2}}$，故 $g^{-1}=a^{-d}\frac{B_0^{2}}{2N_{\mathrm{site}}}\sum'_{\pmb k}(\xi_{\pmb k}^{2}+B_0^{2})^{-3/2}>0$，保证 $(\partial_t\pmb n)^{2}$ 项系数为正。）

\textbf{(5.6.7)：}取 $\hat a=\delta\xi_{\pmb k}\sigma^{y}$：
\begin{equation*}
E^{2}=(\bar\xi_{\pmb k}\pm\delta\xi_{\pmb k})^{2}+B_0^{2}
=\Big[\sqrt{\bar\xi_{\pmb k}^{2}+B_0^{2}}\pm\delta\xi_{\pmb k}\Big]^{2}
\quad\text{（按 }\sqrt{A^{2}+B_0^{2}}\pm\hat a\text{ 重排）},
\end{equation*}
\begin{equation*}
E^{v}_{\pm,\pmb k}=\Big|\sqrt{\tfrac14\big(\xi_{\pmb k+\pmb a/2}+\xi_{\pmb k-\pmb a/2}\big)^{2}+B_0^{2}}\ \pm\ \tfrac12\big(\xi_{\pmb k+\pmb a/2}-\xi_{\pmb k-\pmb a/2}\big)\Big| ,
\end{equation*}
即式 (5.6.7)（等价地 $E^{2}=\xi_{\pmb k\pm\pmb a/2}^{2}+B_0^{2}$：均匀自旋规范势等价于把动量平移 $\pm\pmb a/2$）。同样由 $-\int\mathrm{d}t\,E_0(\pmb a(t))$ 展开，$\pmb a=\partial_{\pmb x}\phi$、$C_2\pmb a^{2}=\int\mathrm{d}^{d}x\,\frac{v^{2}(\partial_{\pmb x}\phi)^{2}}{2g}$，比较得
\begin{equation*}
\boxed{\;\frac{v^{2}}{g}=a^{-d}\frac{\partial^{2}E_0(\pmb a)}{\partial\pmb a^{2}}\Big|_{\pmb a=0}\;}
\end{equation*}
（显式：$E_0(\pmb a)-E_0(0)=-N_{\mathrm{site}}^{-1}\sum'_{\pmb k}\frac{B_0^{2}}{E_{\pmb k}^{3}}\,\delta\xi_{\pmb k}^{2}$，$\delta\xi_{\pmb k}=\frac12a_i\partial_i\xi_{\pmb k}$，故 $v^{2}/g=a^{-d}\frac{B_0^{2}}{4N_{\mathrm{site}}}\sum'_{\pmb k}\frac{(\partial_i\xi_{\pmb k})^{2}}{(\xi_{\pmb k}^{2}+B_0^{2})^{3/2}}$，按方向给出 $v$ 的各向异性。）两式合起来给出 $L_\sigma$ 的全部低能系数，证毕。

% ---------------------------------------------------------------
\section*{问题 5.6.2}
\begin{problem}
从海森堡模型的有效拉格朗日量 (5.6.10) $L=\mathrm{i}\sum_i z_i^{\dagger}\dot z_i-\frac{J}{4}\sum_{\langle ij\rangle}\pmb n_i\cdot\pmb n_j$ 出发，推导式 (5.6.12) 的作用量 $S=\int\mathrm{d}t\,\mathrm{d}x\,\frac{1}{2g}[(\partial_t\pmb n)^{2}-v^{2}(\partial_x\pmb n)^{2}]+\theta W$，即由 $J$ 与晶格间距 $a$ 确定 $g$、$v$（及 $\theta$）。
\end{problem}
\solution

反铁磁涨落取 $\pmb n_i(t)=(-)^{i}\pmb n(x^i,t)+\pmb m(x^i,t)$（$|\pmb n|=1$、$|\pmb m|\ll1$、$\pmb m\cdot\pmb n=0$），其中 $\pmb m$ 是将被积掉的高能铁磁涨落。

\textbf{Berry 项。}由 $\int\mathrm{d}t\,2\mathrm{i}z^{\dagger}\dot z$ 等于 $\pmb n_i$ 所张立体角，正文给出分解
\begin{equation*}
\mathrm{i}\sum_i z_i^{\dagger}\dot z_i
=\int\mathrm{d}x\;\frac14\,\pmb n\cdot(\partial_t\pmb n\times\partial_x\pmb n)
+\int\mathrm{d}x\;\frac{1}{2a}\,\pmb n\cdot(\dot{\pmb n}\times\pmb m) ,
\end{equation*}
（光滑部分：偶、奇两套格点的立体角以同号相加，得 $\frac{1}{2a}\int\mathrm{d}x$ 的 $\frac{1}{2}\pmb n\cdot(\partial_t\pmb n\times\partial_x\pmb n)$ 密度；交错部分产生 $\pmb m$ 的线性耦合）。

\textbf{交换项。}对近邻 $\pmb j=\pmb i\pm a\hat x$：$\pmb n_i\cdot\pmb n_j=-\pmb n\cdot\pmb n'+(-)^{i}(\pmb n\cdot\pmb m'-\pmb m\cdot\pmb n')+\pmb m\cdot\pmb m'$，而 $\pmb n\cdot\pmb n'=1-\frac{a^{2}}{2}(\partial_x\pmb n)^{2}+O(a^{3})$，$\pmb m\cdot\pmb m'=\pmb m^{2}+O(am)$。每格点两条键：
\begin{equation*}
-\frac{J}{4}\sum_{\langle ij\rangle}\pmb n_i\cdot\pmb n_j
=\text{常数}-\frac{Ja}{8}\int\mathrm{d}x\,(\partial_x\pmb n)^{2}
-\frac{J}{2a}\int\mathrm{d}x\,\pmb m^{2}
+(-)^{i}\text{ 型交错项} .
\end{equation*}

\textbf{积掉 $\pmb m$。}拉格朗日量中 $\pmb m$ 的部分为
\begin{equation*}
L_m=\int\mathrm{d}x\,\Big[\frac{1}{2a}\,\pmb n\cdot(\dot{\pmb n}\times\pmb m)-\frac{J}{2a}\pmb m^{2}\Big]
\quad\Longrightarrow\quad
\pmb m=\frac{\pmb n\times\dot{\pmb n}}{2J} ,
\end{equation*}
代回：$\frac{1}{2a}\pmb m\cdot(\pmb n\times\dot{\pmb n})-\frac{J}{2a}\pmb m^{2}=\frac{|\dot{\pmb n}|^{2}}{4aJ}-\frac{|\dot{\pmb n}|^{2}}{8aJ}=\frac{(\partial_t\pmb n)^{2}}{8aJ}$。（交错项给出 $\pmb m$ 与 $\partial_x\pmb n$ 的高阶耦合，属次邻近导数项，舍去。）

\textbf{结果。}
\begin{equation*}
S=\int\mathrm{d}t\,\mathrm{d}x\,\Big[\frac{1}{8aJ}(\partial_t\pmb n)^{2}-\frac{Ja}{8}(\partial_x\pmb n)^{2}\Big]+\theta W ,
\qquad
\frac{\theta}{4\pi}=\frac14\ \Longrightarrow\ \theta=\pi ,
\end{equation*}
与式 (5.6.12) 逐项比较：
\begin{equation*}
\boxed{\;\frac{1}{2g}=\frac{1}{8aJ}\ \Longrightarrow\ g=4aJ;\qquad
\frac{v^{2}}{2g}=\frac{Ja}{8}\ \Longrightarrow\ v^{2}=a^{2}J^{2},\ v=Ja;\qquad \theta=\pi\;}
\end{equation*}
自洽检验：$v=2JSa/\hbar$ 在 $S=\frac12$ 时恰为 $Ja$，与线性自旋波理论的反铁磁自旋波速一致；$\theta=\pi=2\pi S$ 与正文“半整数自旋 $\theta=\pi$”的量子化判据一致（拓扑项不由 $J$、$a$ 这些微扰参数修正，只由平移一格对称性固定为 0 或 $\pi$）。
